\subsection{AMALGAMATED SUM OF MONOIDS}

Let \((M_{i})_{i\in I}\) denote a family of monoids and \(e_i\) the identity element of \(\mathbf{M}_i\) . We are given a monoid \(\mathbf{A}\) and a family of homomorphisms \(h_i\colon \mathbf{A}\to \mathbf{M}_i\) (for \(i\in I\) ).

The set S the sum of the family \((\mathbf{M}_{i})_{i\in\mathbb{I}}\) has elements the ordered pairs \((i,x)\) with \(i\in I\) and \(x\in M_{i}\) . For every triple \(\alpha=(i,x,x^{\prime})\) with \(i\in I\) , \(x,x^{\prime}\) in \(M_{i}\) , write \(u_{\alpha}=(i,xx^{\prime})\) and \(v_{\alpha}=(i,x).(i,x^{\prime})\) ; for every triple \(\lambda=(i,j,a)\) in \(I\times I\times A\) , write \(p_{\lambda}=(i,h_{i}(a))\) and \(q_{\lambda}=(j,h_{j}(a))\) ; for all \(i\in I\) , write \(\varepsilon_{i}=(i,e_{i})\) . The monoid M defined by S and the relators \((u_{\alpha},v_{\alpha})\) , \((p_{\lambda},q_{\lambda})\) and \((\varepsilon_{i},e)\) is called the sum of the family \((\mathbf{M}_{i})_{i\in\mathbb{I}}\) amalgamated by A. Let \(\phi\) denote the canonical homomorphism of \(\mathrm{Mo}(\mathrm{S})\) onto M and write \(\phi_{i}(x)=\phi(i,x)\) for \((i,x)\in\mathrm{S}\) . It is said that \(\phi_{i}\) is the canonical mapping of \(M_{i}\) into M. For all \(a\in A\) , the element \(\phi(i,h_{i}(a))\) is independent of i and denoted by \(h(a).†\)

The universal property of monoids defined by generators and relators (no. 2) implies the following result:

PROPOSITION 4. (a) For all \(i \in \mathbf{I}\) , the mapping \(\phi_i\) is a homomorphism of \(\mathbf{M}_i\) into \(\mathbf{M}\) and \(\phi_i \circ h_i = h\) for all \(i \in \mathbf{I}\) . Further, \(\mathbf{M}\) is generated by \(\bigcup_{i \in \mathbf{I}} \phi_i(\mathbf{M}_i)\) .

\begin{enumerate}
\def\labelenumi{(\alph{enumi})}
\setcounter{enumi}{1}
\tightlist
\item
  Let \(\mathbf{M}'\) be a monoid and \(f': \mathbf{M}_i \to \mathbf{M}'\) (for \(i \in \mathbf{I}\) ) homomorphisms such that \(f_i \circ h_i\) is independent of \(i \in \mathbf{I}\) . There exists one and only one homomorphism \(f: \mathbf{M} \to \mathbf{M}'\) such that \(f_i = f \circ \phi_i\) for all \(i \in \mathbf{I}\) .
\end{enumerate}

In what follows we shall make the following hypothesis:

\begin{enumerate}
\def\labelenumi{(\Alph{enumi})}
\tightlist
\item
  For all \(i \in \mathbf{I}\) , there exists a subset \(\mathbf{P}_i\) of \(\mathbf{M}_i\) containing \(e_i\) such that the mapping \((a, p) \mapsto h_i(a).p\) of \(\mathbf{A} \times \mathbf{P}_i\) into \(\mathbf{M}_i\) is bijective.
\end{enumerate}

It implies that the homomorphisms \(h_i\) are injective. Let \(x \in \mathbf{M}\) ; every finite sequence \(\sigma = (a: i_1, \ldots, i_n; p_1, \ldots, p_n)\) with \(a \in \mathbf{A}\) , \(i_\alpha \in \mathbf{I}\) and \(p_\alpha \in \mathbf{P}_{i_\alpha}\) for \(1 \leqslant \alpha \leqslant n\) , satisfying

\[
x = h (a) \cdot \prod_ {\alpha = 1} ^ {n} \phi_ {i _ {\alpha}} (p _ {\alpha})\tag{3}
\]

is called a decomposition of \(x\) . The integer \(n \geqslant 0\) is called the length of the decomposition \(\sigma\) and is denoted by \(l(\sigma)\) ; the sequence \((e)\) is a decomposition of length 0 of the identity element of \(M\) . The decomposition \(\sigma\) is called reduced if \(i_{\alpha} \neq i_{\alpha + 1}\) for \(1 \leqslant \alpha < n\) and \(p_{\alpha} \neq e_{i_{\alpha}}\) for \(1 \leqslant \alpha \leqslant n\) .

PROPOSITION 5. Under hypothesis (A) every element \(x\) of \(\mathbf{M}\) admits a unique reduced decomposition \(\sigma\) . Every decomposition \(\sigma' \neq \sigma\) of \(x\) satisfies \(l(\sigma') > l(\sigma)\) .

\BourbakiRunInHeading{(A) Uniqueness of a reduced decomposition:}

Let \(\Sigma\) denote the set of sequences \(\sigma = (a; i_1, \ldots, i_n; p_1, \ldots, p_n)\) with \(n \geqslant 0\) , \(a \in \mathbf{A}\) , \(i_\alpha \in \mathbf{I}\) and \(p_\alpha \in \mathrm{P}_{i_\alpha} - \{e_{i_\alpha}\}\) for \(1 \leqslant \alpha \leqslant n\) , such that \(i_\alpha \neq i_{\alpha+1}\) for \(1 \leqslant \alpha < n\) . Let \(\Phi\) denote the mapping of \(\Sigma\) into \(\mathbf{M}\) defined by

\[
\Phi (a; i _ {1}, \dots , i _ {n}; p _ {1}, \dots , p _ {n}) = h (a). \prod_ {\alpha = 1} ^ {n} \phi_ {i _ {\alpha}} (p _ {\alpha}).\tag{4}
\]

A reduced decomposition of \(x \in \mathbf{M}\) is an element \(\sigma\) of \(\Sigma\) such that \(\Phi(\sigma) = x\) . For all \(i \in \mathbf{I}\) , let \(\Sigma_i\) be the subset of \(\Sigma\) consisting of the sequences \((e; i_1, \ldots, i_n; p_1, \ldots, p_n)\) with \(i \neq i_1\) when \(n > 0\) . Let

\[
\sigma = (e; i _ {1}, \dots , i _ {n}; p _ {1}, \dots , p _ {n})
\]

be in \(\Sigma_{i}\) and \(\xi\) in \(\mathbf{M}_i\) ; let \(\xi = h_i(a).p\) with \(a\in \mathbf{A}\) and \(p\in \mathbf{P}_i\) , and

\[
\Psi_ {i} (\xi , \sigma) = \left\{ \begin{array}{l l} (a; i _ {1}, \dots , i _ {n}; p _ {1}, \dots , p _ {n}) & \text { if } p = e _ {i} \\ (a; i, i _ {1}, \dots , i _ {n}; p, p _ {1}, \dots , p _ {n}) & \text { if } p \neq e _ {i}. \end{array} \right.\tag{5}
\]

It is immediate that \(\Psi_{i}\) is a bijection of \(\mathbf{M}_i\times \Sigma_i\) onto \(\Sigma\) .

Let \(i \in I\) and \(x \in M_i\) ; as \(\Psi_i\) is bijective, a mapping \(f_{i,x}\) of \(\Sigma\) into itself is defined by

\[
f _ {i, x} (\Psi_ {i} (\xi , \sigma)) = \Psi_ {i} (x \xi , \sigma) \quad (\xi \in \mathbf {M} _ {i}, \sigma \in \Sigma_ {i}).\tag{6}
\]

Further, for \(a \in \mathbf{A}, f_a\) denotes the mapping of \(\Sigma\) into itself defined by

\[
f _ {a} (a ^ {\prime}; i _ {1}, \dots , i _ {n}; p _ {1}, \dots , p _ {n}) = (a a ^ {\prime}; i _ {1}, \dots , i _ {n}; p _ {1}, \dots , p _ {n}).\tag{7}
\]

Clearly \(f_{i, e_i}\) is the identity mapping of \(\Sigma\) and \(f_{i, xx'} = f_{i, x} \circ f_{i, x'}\) for \(x, x'\) in \(M_i\) and \(f_{i, h_i(a)} = f_a\) for \(a \in A\) and \(i \in I\) .

Then Proposition 4 may be applied to the case where \(\mathbf{M}'\) is the monoid of mappings of \(\Sigma\) into itself with law of composition \((f, f') \mapsto f \circ f'\) and where \(f_{i}\) is the homomorphism \(x \mapsto f_{i,x}\) of \(\mathbf{M}_{i}\) into \(\mathbf{M}'\) ; then there exists a homomorphism \(f\) of \(\mathbf{M}\) into \(\mathbf{M}'\) such that \(f_{i,x} = f(\phi_i(x))\) for \(i \in \mathbf{I}\) and \(x \in \mathbf{M}_i\) . Let

\[
\sigma = (a; i _ {1}, \dots , i _ {n}; p _ {1}, \dots , p _ {n})
\]

be in \(\Sigma\) . Formulae (5) to (7) imply by induction on \(n\) the relation

\[
\begin{array}{l} \sigma = (f _ {a} \circ f _ {i _ {1}, p _ {1}} \circ \dots \circ f _ {i _ {n}, p _ {n}}) (e) \\ = f (h (a) \phi_ {i _ {1}} (p _ {1}) \dots \phi_ {i _ {n}} (p _ {n})) (e), \end{array}
\]

that is \(\sigma = f(\Phi(\sigma))(e)\) . This proves that \(\Phi\) is injective.

\BourbakiRunInHeading{(B) Existence of a decomposition:}

Let D be the set of elements of M admitting a decomposition. Then \(e \in D\) and M is generated by \(\bigcup_{i \in I} \phi_i(M_i)\) and hence by \(h(A) \cup \bigcup_{i \in I} \phi_i(P_i)\) . Then \(D \cdot \phi_i(P_i) \subset D\) for all \(i \in I\) ; to prove that D = M, it thus suffices to prove the relation \(D \cdot h(A) \subset D\) . This follows from the following more precise lemma:

Lemma 1. Let \(i_1, \ldots, i_n\) be in I and \(p_\alpha\) in \(\mathbf{P}_{i_\alpha}\) for \(1 \leqslant \alpha \leqslant n\) . For all \(a \in A\) there exists \(a' \in A\) and a sequence \((p'_\alpha)_{1 \leqslant \alpha \leqslant n}\) with \(p'_\alpha \in \mathbf{P}_{i_\alpha}\) such that

\[
\phi_ {i _ {1}} (p _ {1}) \dots \phi_ {i _ {n}} (p _ {n}) h (a) = h (a ^ {\prime}) \phi_ {i _ {1}} (p _ {1} ^ {\prime}) \dots \phi_ {i _ {n}} (p _ {n} ^ {\prime}).
\]

\(h(a) = \phi_{i_n}(h_{i_n}(a))\) and there exists \(a_{n} \in \mathbf{A}\) and \(p_n' \in \mathbf{P}_{i_n}\) with

\[
p _ {n} \cdot h _ {i _ {n}} (a) = h _ {i _ {n}} (a _ {n}) \cdot p _ {n} ^ {\prime}.
\]

It follows that \(\phi_{i_n}(p_n)h(a) = h(a_n)\phi_{i_n}(p_n')\) , whence

\[
\phi_ {i _ {1}} (p _ {1}) \dots \phi_ {i _ {n - 1}} (p _ {n - 1}) \phi_ {i _ {n}} (p _ {n}) h (a) = \phi_ {i _ {1}} (p _ {1}) \dots \phi_ {i _ {n - 1}} (p _ {n - 1}) h (a _ {n}) \phi_ {i _ {n}} (p _ {n} ^ {\prime});
\]

the lemma follows from this by induction on n.

\BourbakiRunInHeading{(C) End of the proof:}

Let \(x \in M\) and let n be the minimum of the lengths of decompositions of x. We shall prove that every decomposition \(\sigma\) of x of length n is reduced. This will establish the existence of a reduced decomposition of x; the uniqueness of the reduced decomposition then implies \(l(\sigma') > l(\sigma)\) for every decomposition \(\sigma' \neq \sigma\) of x.

The case \(n = 0\) being trivial, suppose \(n > 0\) . Let

\[
\sigma = (a; i _ {1}, \dots , i _ {n}; p _ {1}, \dots , p _ {n})
\]

be a decomposition of \(x\) of length \(n\) . If there existed an integer \(\alpha\) with \(1 \leqslant \alpha \leqslant n\) and \(p_{\alpha} = e_{i_{\alpha}}\) , the sequence

\[
(a; i _ {1}, \dots , i _ {\alpha - 1}, i _ {\alpha + 1}, \dots , i _ {n}; p _ {1}, \dots , p _ {\alpha - 1}, p _ {\alpha + 1}, \dots , p _ {n})
\]

would be a decomposition of \(x\) of length \(n - 1\) , which is excluded. Suppose that there exists an integer \(\alpha\) with \(1 \leqslant \alpha < n\) and \(i_{\alpha} = i_{\alpha + 1}\) and let

\[
p _ {\alpha} p _ {\alpha + 1} = h _ {i _ {\alpha}} \left(a ^ {\prime}\right) \cdot p _ {\alpha} ^ {\prime}
\]

with \(a' \in \mathbf{A}\) and \(p_{\alpha}' \in \mathrm{P}_{t_{\alpha}}\) ; by Lemma 1 there exists elements \(a'' \in \mathbf{A}\) ,

\[
p _ {1} ^ {\prime} \in \mathrm{P} _ {i _ {1}}, \dots , p _ {\alpha - 1} ^ {\prime} \in \mathrm{P} _ {i _ {\alpha - 1}}
\]

such that

\[
\phi_ {i _ {1}} (p _ {1}) \dots \phi_ {i _ {\alpha - 1}} (p _ {\alpha - 1}) h (a ^ {\prime}) = h (a ^ {\prime \prime}) \phi_ {i _ {1}} (p _ {1} ^ {\prime}) \dots \phi_ {i _ {\alpha - 1}} (p _ {\alpha - 1} ^ {\prime})
\]

and the sequence

\[
\left(a a ^ {\prime \prime}; i _ {1}, \dots , i _ {\alpha - 1}, i _ {\alpha}, i _ {\alpha + 2}, \dots , i _ {n}; p _ {1} ^ {\prime}, \dots , p _ {\alpha - 1} ^ {\prime}, p _ {\alpha} ^ {\prime}, p _ {\alpha + 2}, \dots , p _ {n}\right)
\]

is a decomposition of \(x\) of length \(n - 1\) , which is a contradiction.

We have thus proved that \(\sigma\) is reduced.

COROLLARY. Under hypothesis (A) the homomorphisms \(\phi_i\) and \(h\) are injective. For \(i \neq j\) in I, \(\phi_i(M_i) \cap \phi_j(M_j) = h(A)\) .

First \(h\) is injective: if \(h(a) = h(a')\) , then (a) and \((a')\) are two reduced decompositions of the same element of \(\mathbf{M}\) , whence \(a = a'\) . Let \(i \in \mathbf{I}\) ; then \(h(\mathbf{A}) = \phi_i(h_i(\mathbf{A})) \subset \phi_i(\mathbf{M}_i)\) ; the uniqueness of reduced decompositions implies

\[
h (\mathrm{A}) \cap \phi_ {i} (\mathrm{M} _ {i} - h _ {i} (\mathrm{A})) = \varnothing ,
\]

whence \(\phi_{i}(\mathbf{M}_{i} - h_{i}(\mathbf{A})) = \phi_{i}(\mathbf{M}_{i}) - h(\mathbf{A}).\)

The injectivity of the homomorphisms \(\phi_{i}\) and the relation

\[
\phi_ {i} (\mathbf {M} _ {i}) \cap \phi_ {j} (\mathbf {M} _ {j}) \subset h (\mathrm{A})
\]

for \(i \neq j\) are then consequences of the following fact: for \(i, j\) in \(I, x\) in \(M_i = h_i(A)\) and \(y\) in \(M_j = h_j(A)\) , the relation \(\phi_i(x) = \phi_j(y)\) implies \(i = j\) and \(x = y\) . Let \(x = h_i(a).p\) and \(y = h_j(b).q\) with \(a, b\) in \(A, p\) in \(P_i = \{e_i\}\) and \(y\) in \(P_j = \{e_j\}\) . Then \(\phi_i(x) = h(a)\phi_i(p)\) and \(\phi_j(y) = h(b)\phi_j(q)\) and hence \((a; i; p)\) and \((b; j; q)\) are two reduced decompositions of the same element of \(M\) . It follows that \(i = j, a = b\) and \(p = q\) , whence \(x = h_i(a)p = h_j(b)q = y\) .

When hypothesis (A) is fulfilled, we shall identify each monoid \(M_{i}\) with a submonoid of M by means of \(\phi_{i}\) ; similarly, we shall identify A with a submonoid of M by h. Then M is generated by \(\bigcup_{i\in I}M_{i}\) and \(M_{i}\cap M_{j}=A\) for \(i\neq j\) .

Every element of M can be written uniquely in the form \(a.\prod_{\alpha=1}p_{\alpha}\) with \(a\in A, p_{1}\in P_{i1}=\{e\},\ldots,p_{n}\in P_{in}=\{e\}\) and \(i_{\alpha}\neq i_{\alpha+1}\) for \(1\leqslant\alpha<n\) . Finally, if \(M'\) is a monoid and \((f_{i}:M_{i}\to M')\) (for \(i\in I\) ) a family of homomorphisms whose restrictions to A are the same homomorphism of A into \(M'\) , there exists one and only one homomorphism \(f:M\to M'\) inducing \(f_{i}\) on \(M_{i}\) for all \(i\in I\) .

Hypothesis (A) is satisfied in two important cases:

\begin{enumerate}
\def\labelenumi{(\alph{enumi})}
\tightlist
\item
  \(\mathbf{A} = \{e\}\) . In this case, there is a family \((\mathbf{M}_i)_{i\in I}\) of monoids and \(\mathbf{M}\) is called the monoidal sum of this family. Each \(\mathbf{M}_i\) is identified with a submonoid of \(\mathbf{M}\) and \(\mathbf{M}\) is generated by \(\bigcup_{i\in I}\mathbf{M}_i\) ; further, \(\mathbf{M}_i\cap \mathbf{M}_j = \{e\}\) for \(i\neq j\) . Every element of \(\mathbf{M}\) may be written uniquely in the form \(x_{1}\ldots x_{n}\) with
\end{enumerate}

\[
x _ {1} \in \mathrm{M} _ {i _ {1}} - \{e \}, \dots , x _ {n} \in \mathrm{M} _ {i _ {n}} - \{e \}
\]

and \(i_{\alpha} \neq i_{\alpha + 1}\) for \(1 \leqslant \alpha \leqslant n\) . Finally, for every family of homomorphisms \((f_i: \mathbf{M}_i \to \mathbf{M}')\) , there exists a unique homomorphism \(f: \mathbf{M} \to \mathbf{M}'\) whose restriction to \(\mathbf{M}_i\) is \(f_i\) for all \(i \in I\) .

\begin{enumerate}
\def\labelenumi{(\alph{enumi})}
\setcounter{enumi}{1}
\tightlist
\item
  There is a family of groups \((\mathbf{G}_i)_{i\in I}\) containing as subgroup the same group A and \(h_i\) is the injection of A into \(\mathbf{G}_i\) . The sum of the family \((\mathbf{G}_i)_{i\in I}\) amalgamated by A is then a group G: the monoid G is generated by \(\bigcup_{i\in I}\phi_i(\mathbf{G}_i)\) and every element of \(\bigcup_{i\in I}\phi_i(\mathbf{G}_i)\) admits an inverse in G (cf.~§ 2, no. 3, Corollary 1 to Proposition 4); it is denoted by \(\star_{\mathrm{A}}\mathrm{G}_{i}\) or \(\mathrm{G}_{1}*\mathrm{A}\mathrm{G}_{2}\) when \(\mathbf{I} = \{1,2\}\) . When A consists of the identity element, it is also said that G is the free product of the family \((\mathbf{G}_i)_{i\in I}\) of groups and it is denoted by \(\star \mathbf{G}_i\) (or \(\mathbf{G}_1*\mathbf{G}_2\) if \(\mathbf{I} = \{1,2\}\) ).†
\end{enumerate}

\subsection{APPLICATION TO FREE MONOIDS}

Lemma 2. Let \(\mathbf{M}\) be the monoidal sum of the family \((M_x)_{x\in \mathbf{X}}\) defined by \(\mathbf{M}_x = \mathbf{N}\) for all \(x\in \mathbf{X}\) and let \(\phi_{x}\) denote the canonical homomorphism of \(\mathbf{M}_x\) into \(\mathbf{M}\) . The mapping \(x\mapsto \phi_x(1)\) of \(\mathbf{X}\) into \(\mathbf{M}\) extends to an isomorphism \(h\) of \(\operatorname {Mo}(\mathbf{X})\) onto \(\mathbf{M}\) .

Let \(h\) be the homomorphism of \(\operatorname{Mo}(\mathbf{X})\) into \(\mathbf{M}\) characterized by \(h(x) = \phi_x(1)\) . For every integer \(n \geqslant 0\) , \(\phi_x(n) = \phi_x(1)^n = h(x)^n\) and as \(\mathbf{M}\) is generated by \(\bigcup_{x \in \mathbf{X}} \phi_x(\mathbf{N})\) , it is also generated by \(h(\mathbf{X})\) . Hence \(h\) is surjective. Moreover, for all \(x\) in \(\mathbf{X}\) , the mapping \(n \mapsto x^n\) is a homomorphism of \(\mathbf{N} = \mathbf{M}_x\) into \(\operatorname{Mo}(\mathbf{X})\) ; there thus exists (no. 3, Proposition 4) a homomorphism \(h'\) of \(\mathbf{M}\) into \(\operatorname{Mo}(\mathbf{X})\) such that \(h'(\phi_x(n)) = x^n\) for \(x \in \mathbf{X}\) and \(n \in \mathbf{N}\) ; in particular, \(h'(h(x)) = x\) for \(x \in \mathbf{X}\) and hence \(h' \circ h\) is the identity homomorphism of \(\operatorname{Mo}(\mathbf{X})\) . Therefore \(h\) is injective. It has thus been proved that \(h\) is bijective.

PROPOSITION 6. Let \(w\) be an element of \(\mathbf{Mo}(\mathbf{X})\) .

\begin{enumerate}
\def\labelenumi{(\alph{enumi})}
\tightlist
\item
  There exist an integer \(n \geqslant 0\) , elements \(x_{\alpha}\) of \(\mathbf{X}\) and integers \(m(\alpha) > 0\) (for
\end{enumerate}

\(1 \leqslant \alpha \leqslant n\) such that \(x_{\alpha} \neq x_{\alpha + 1}\) for \(1 \leqslant \alpha < n\) and \(w = \prod_{\alpha = 1}^{n} x_{\alpha}^{m(\alpha)}\) . The sequence \((x_{\alpha}, m(\alpha))_{1 \leqslant \alpha \leqslant n}\) is determined uniquely by these conditions.

\[
x _ {\beta} ^ {\prime}
\]

\[
m ^ {\prime} (\beta)
\]

\[
1 \leqslant \beta \leqslant p
\]

(b) \(w = \prod_{\beta=1}^{p} x'^{m'(\beta)}.\) Then \(p \geqslant n\) . If \(p = n\) , then \(x_{\beta}' = x_{\beta}\) and \(m'(\beta) = m(\beta)\) for \(1 \leqslant \beta \leqslant p\) .

In the notation of Lemma 2, \(\bar{h}^{1}(\phi_x(n)) = x^n\) for \(x \in \mathbf{X}\) and \(n \in \mathbf{N}\) . Proposition 6 then follows from no. 3, Proposition 5.

\subsection{FREE GROUPS}

Let \(G_x = Z\) for all \(x \in X\) . The free product of the family \((G_x)_{x \in X}\) is called the free group constructed on \(X\) and is denoted by \(F(X)\) . Let \(\phi_x\) denote the canonical homomorphism of \(G_x = Z\) into \(F(X)\) . By no. 3, Corollary to Proposition 5, the mapping \(x \mapsto \phi_x(1)\) of \(X\) into \(F(X)\) is injective; we shall identify \(X\) with its image in \(F(X)\) under this mapping. Then \(X\) generates \(F(X)\) and \(e \notin X\) .

Applying no. 3, Proposition 5, we obtain the following result:

PROPOSITION 7. Let \(g\) be an element of the free group \(F(X)\) . There exist an integer \(n \geqslant 0\) and a sequence \((x_{\alpha}, m(\alpha))_{1 \leqslant \alpha \leqslant n}\) determined uniquely by the relations \(x_{\alpha} \in X\) ,

\[
g = \prod_ {\alpha = 1} ^ {n} x _ {\alpha} ^ {m (\alpha)}
\]

The free group \(F(X)\) enjoys the following universal property:

PROPOSITION 8. Let G be a group and f a mapping of X into G. There exists one and only one homomorphism \(\bar{f}\) of F(X) into G which extends \(f\) .

The uniqueness of \(\bar{f}\) follows from the fact that the group \(\mathrm{F}(\mathbf{X})\) is generated by \(\mathbf{X}\) . For all \(x\) in \(\mathbf{X}\) , let \(f_x\) be the homomorphism \(n \mapsto f(x)^n\) of \(\mathbf{Z}\) into \(\mathbf{G}\) . By no. 3, Proposition 4, there exists a homomorphism \(\bar{f}\) of \(\mathrm{F}(\mathbf{X})\) into \(\mathbf{G}\) such that \(\bar{f}(x^n) = f_x(n)\) for \(x \in \mathbf{X}\) and \(n \in \mathbf{Z}\) ; in particular, \(\bar{f}(x) = f_x(1) = f(x)\) for all \(x \in \mathbf{X}\) and hence \(\bar{f}\) extends \(f\) .

Let \(u: X \to Y\) be a mapping. By Proposition 8 there exists one and only one homomorphism of \(F(X)\) into \(F(Y)\) which coincides with \(u\) on \(X\) ; it is denoted by \(F(u)\) . As in the case of magmas (no. 1) the formula

\[
\mathbf {F} (v \circ u) = \mathbf {F} (v) \circ \mathbf {F} (u)
\]

is established for every mapping \(v: \mathbf{Y} \to \mathbf{Z}\) and it is shown that \(\mathrm{F}(u)\) is injective (resp. surjective, bijective) if \(u\) is. For every subset \(S\) of \(X\) , \(F(S)\) will be identified with the subgroup of \(F(X)\) generated by \(S\) .

Let I be a set. In certain cases it is of interest not to identify i in I with its canonical image \(\phi_i(1)\) in the free group \(\mathbf{F}(\mathbf{I})\) ; the latter will be denoted by \(\mathbf{T}_i\) (or \(\mathbf{T}_i', \mathbf{X}_i, \ldots\) as the case may be) and called the indeterminate of index \(i\) . The free group \(\mathbf{F}(\mathbf{I})\) is then denoted by \(\mathbf{F}((\mathbf{T}_i)_{i \in 1})\) or \(\mathbf{F}(\mathbf{T}_1, \ldots, \mathbf{T}_n)\) if \(\mathbf{I} = \{1, 2, \ldots, n\}\) .

Let G be a group and \(\mathbf{t} = (t_{i})_{i \in \mathbf{I}}\) a family of elements of G. By Proposition 8 there exists a homomorphism \(f_{t}\) of \(\mathrm{F}((\mathrm{T}_{i})_{i \in \mathbf{I}})\) into G characterized by \(f_{\mathbf{t}}(\mathrm{T}_{i}) = t_{i}\) for all \(i \in I\) . The image of an element w of \(\mathrm{F}((\mathrm{T}_{i})_{i \in \mathbf{I}})\) under \(f_{t}\) will be denoted by \(w(\mathbf{t})\) or \(w(t_{1}, \ldots, t_{n})\) if \(I = \{1, 2, \ldots, n\}\) ; \(w(\mathbf{t})\) is said to result from the substitution \(T_{i} \mapsto t_{i}\) in w. In particular, if we take \(\mathrm{G} = \mathrm{F}((\mathrm{T}_{i})_{i \in \mathbf{I}})\) and \((t_{i}) = (\mathrm{T}_{i}) = \mathbf{T}\) , \(f_{T}\) is the identity homomorphism of G, whence \(w(\mathbf{T}) = w\) ; for \(I = \{1, 2, \ldots, n\}\) , then \(w(\mathrm{T}_{1}, \ldots, \mathrm{T}_{n}) = w\) .

Let G and G' be two groups, u a homomorphism of G into G' and \(\mathbf{t} = (t_{1}, \ldots, t_{n})\) a finite sequence of elements in G. Let \(\mathbf{t}' = (u(t_{1}), \ldots, u(t_{n}))\) ; the homomorphism \(u \circ f_{t}\) of \(\mathrm{F}(\mathrm{T}_{1}, \ldots, \mathrm{T}_{n})\) into G' maps \(T_{i}\) to \(u(t_{i})\) for \(1 \leqslant i \leqslant n\) and hence is equal to \(f_{t'}\) ; for w in \(\mathrm{F}(\mathrm{T}_{1}, \ldots, \mathrm{T}_{n})\) , then

\[
u (w (t _ {1}, \dots , t _ {n})) = w (u (t _ {1}), \dots , u (t _ {n})).\tag{8}
\]

Let \(w\) be given in \(F(T_1, \ldots, T_n)\) and elements \(v_1, \ldots, v_n\) in the free group \(F(T_1', \ldots, T_m')\) . The substitution \(T_i \mapsto v_i\) defines an element \(w' = w(v_1, \ldots, v_n)\) of \(F(T_1', \ldots, T_m')\) . Let \(G\) be a group, \(t_1, \ldots, t_m\) elements of \(G\) and \(u\) the homomorphism of \(F(T_1', \ldots, T_m')\) into \(G\) characterized by \(u(T_j') = t_j\) for \(1 \leqslant j \leqslant m\) . Then \(u(v_i) = v_i(t_1, \ldots, t_m)\) and \(u(w') = w(t_1, \ldots, t_m)\) ; formula (8) thus implies

\[
w ^ {\prime} (t _ {1}, \dots , t _ {m}) = w (v _ {1} (t _ {1}, \dots , t _ {m}), \dots , v _ {n} (t _ {1}, \dots , t _ {m})).\tag{9}
\]

This justifies the ``functional notation'' \(w(t_{1},\ldots,t_{n})\) . The reader is left to extend formulae (8) and (9) to the case of arbitrary indexing sets.

\subsection{PRESENTATIONS OF A GROUP}

Let G be a group and \(\mathbf{t} = (t_{i})_{i \in I}\) a family of elements of G. Let \(f_{t}\) be the unique homomorphism of the free group F(I) into G which maps i to \(t_{i}\) . The image of \(f_{t}\) is the subgroup generated by the elements \(t_{i}\) of G. The elements of the kernel of \(f_{t}\) are called the relators of the family t. t is called generating (resp. free, basic) if \(f_{t}\) is surjective (resp. injective, bijective).

Let G be a group. A presentation of G is an ordered pair \((\mathbf{t}, \mathbf{r})\) consisting of a generating family \(\mathbf{t} = (t_{i})_{i \in I}\) and a family \(\mathbf{r} = (r_{j})_{j \in J}\) of relators such that the kernel \(N_{t}\) of \(f_{t}\) is generated by the elements \(gr_{j}g^{-1}\) for \(g \in \mathrm{F}(\mathrm{I})\) and \(j \in J\) . It amounts to the same to say that \(N_{t}\) is the normal subgroup of F(I) generated by the \(r_{j}\) for \(j \in J\) (in other words, the smallest normal subgroup of F(I) containing the elements \(r_{j}\) ( \(j \in J\) ), cf.~§4, no. 4). By an abuse of language the generators \(t_{i}\) and the relations \(r_{j}(\mathbf{t}) = e\) are said to constitute a presentation of the group G.

Let I be a set and \(\mathbf{r} = (r_j)_{j \ensuremath{\in} J}\) a family of elements of the free group F(I). Let N(r) be the normal subgroup of F(I) generated by the \(r_j\) for \(j \ensuremath{\in} J\) . Let F(I, r) = F(I)/N(r) and \(\tau_i\) denote the class of \(i\) modulo N(r). The ordered pair \((\tau, \mathbf{r})\) with \(\tau = (\tau_i)_{i \ensuremath{\in} I}\) is a presentation of the group F(I, r); if G is a group and (t, r) is a presentation of G with t = (t\_i)\_\{i \ensuremath{\in} I\}, there exists a unique isomorphism \(u\) of F(I, r) onto G such that \(u(\tau_i) = t_i\) for all \(i \ensuremath{\in} I\) . The group F(I, r) is said to be defined by the generators \(\tau_i\) and the relators \(r_j\) , or by an abuse of language that it is defined by the generators \(\tau_i\) and the relations \(r_j(\tau) = e\) . When I = {[}1, n{]} and J = {[}1, m{]}, it is said that F(I, r) is defined by the presentation

\[
\langle \tau_ {1}, \dots , \tau_ {n}; r _ {1}, \dots , r _ {m} \rangle .
\]

If \(r_j = u_j^{-1}v_j\) with \(u_j\) and \(v_j\) in \(\mathrm{F(I)}\) , this presentation is equally denoted by the symbol

\[
\langle \tau_ {1}, \dots , \tau_ {n}; u _ {1} = v _ {1}, \dots , u _ {m} = v _ {m} \rangle .
\]

Examples. (1) The group defined by the presentation \(\langle \tau; \tau^q = e \rangle\) is cyclic of order \(q\) .

\begin{enumerate}
\def\labelenumi{(\arabic{enumi})}
\setcounter{enumi}{1}
\tightlist
\item
  The group defined by the presentation \(\langle x, y; xy = yx \rangle\) is isomorphic to \(\mathbf{Z} \times \mathbf{Z}\) .
\end{enumerate}

PROPOSITION 9. Let G be a group, \(\mathbf{t} = (t_i)_{i \in I}\) a generating family of G and \(\mathbf{r} = (r_j)_{j \in J}\) a family of relators of \(\mathbf{t}\) . The following conditions are equivalent:

\begin{enumerate}
\def\labelenumi{(\alph{enumi})}
\item
  The ordered pair \((\mathbf{t},\mathbf{r})\) is a presentation of G.
\item
  Let \(\mathbf{G}'\) be a group and \(\mathbf{t}' = (t_i')_{i \in \mathbf{I}}\) a family of elements of \(\mathbf{G}'\) . If \(r_j(\mathbf{t}') = e\) for all \(j \in \mathbf{J}\) , there exists a homomorphism \(u\) of \(\mathbf{G}\) into \(\mathbf{G}'\) such that \(u(t_i) = t_i'\) for all \(i \in \mathbf{I}\) .
\item
  Let \(\overline{\mathbf{G}}\) be a group and \(\overline{\mathbf{t}} = (\bar{t}_i)_{i\in \mathbf{I}}\) a generating family of \(\overline{\mathbf{G}}\) such that \(r_j(\overline{\mathbf{t}}) = e\) for all \(j\in \mathbf{J}\) . Every homomorphism of \(\overline{\mathbf{G}}\) into \(\mathbf{G}\) which maps \(\bar{t}_i\) to \(t_i\) for all \(i\in \mathbf{I}\) is an isomorphism.
\end{enumerate}

Let \(f\) denote the homomorphism of \(\mathbf{F}(\mathbf{I})\) into \(\mathbf{G}\) which maps \(i\) to \(t_i\) for all \(i \in \mathbf{I}\) and \(\mathbf{N}\) the kernel of \(f\) .

\begin{enumerate}
\def\labelenumi{(\alph{enumi})}
\item
  \(\Rightarrow\) (b): Suppose that \((\mathbf{t},\mathbf{r})\) is a presentation of G and let \(\mathbf{t}' = (t_i')_{i\in I}\) be a family of elements of a group \(\mathbf{G}'\) with \(r_j(\mathbf{t}') = e\) for all \(j\in J\) . Let \(f'\) be the homomorphism of F(I) into \(\mathbf{G}'\) characterized by \(f'(i) = t_i'\) for all \(i\in I\) . By hypothesis \(f'(r_j) = e\) for all \(j\in J\) and, as N is generated by the elements \(gr_jg^{-1}\) for \(j\in J\) and \(g\in \mathrm{F}(\mathrm{I}),f'(\mathrm{N}) = \{e\}\) . As the homomorphism \(f\colon \mathrm{F}(\mathrm{I})\to \mathrm{G}\) is surjective with kernel N, there exists a homomorphism \(u\colon \mathrm{G}\to \mathrm{G}'\) such that \(f' = u\circ f\) . Then \(u(t_i) = u(f(i)) = f'(i) = t_i'\) .
\item
  \(\Rightarrow\) (c): Suppose condition (b) holds. Let \(\mathbf{t} = (t_i)_{i\in I}\) be a generating family of a group G such that \(r_j(\mathbf{t}) = e\) for all \(j\in J\) and let \(v\) be a homomorphism of \(\overline{\mathbf{G}}\) into G such that \(v(t_i) = t_i\) for all \(i\in I\) . As the family \((t_i)_{i\in I}\) generates G, the homomorphism \(v\) is surjective. By property (b) there exists a homomorphism \(u\colon G\to \overline{\mathbf{G}}\) such that \(u(t_i) = \bar{t}_i\) for all \(i\in I\) . Then \(u(v(\bar{t}_i)) = \bar{t}_i\) for all \(i\in I\) and hence \(u \circ v\) is the identity on \(\overline{G}\) , which proves that \(v\) is injective. Hence \(v\) is an isomorphism and condition (c) holds.
\item
  \(\Rightarrow\) (a): Suppose condition (c) holds. Let \(t_i'\) be the canonical image of \(i\) in \(\mathbf{F}(\mathbf{I},\mathbf{r})\) and \(\mathbf{t}' = (t_i')_{i\in \mathbf{I}}\) ; then \(r_j(\mathbf{t}') = e\) for all \(j\in \mathbf{J}\) . As \(r_j(\mathbf{t}) = e\) for all \(j\in \mathbf{J}\) , there exists one and only one homomorphism \(v\) of \(\mathbf{F}(\mathbf{I},\mathbf{r})\) into \(\mathbf{G}\) such that \(v(\bar{t}_i) = t_i\) for all \(i\in \mathbf{I}\) . By (c), \(v\) is an isomorphism of \(\mathbf{F}(\mathbf{I},\mathbf{r})\) onto \(\mathbf{G}\) which transforms the presentation \((\mathbf{t}',\mathbf{r})\) of \(\mathbf{F}(\mathbf{I},\mathbf{r})\) into a presentation \((\mathbf{t},\mathbf{r})\) of \(\mathbf{G}\) .
\end{enumerate}

\subsection{FREE COMMUTATIVE GROUPS AND MONOIDS}

The set \(\mathbf{Z}^{\mathbf{x}}\) of all mappings of \(\mathbf{X}\) into \(\mathbf{Z}\) is a commutative group under the law defined by \((\alpha + \beta)(x) = \alpha(x) + \beta(x)(\alpha, \beta \in \mathbf{Z}^{\mathbf{x}}, x \in \mathbf{X})\) ; the elements of \(\mathbf{Z}^{\mathbf{x}}\) are sometimes called multiindices. The identity element, denoted by 0, is the constant mapping with value 0. For \(\alpha \in \mathbf{Z}^{\mathbf{x}}\) , the set \(S_{\alpha}\) of \(x \in \mathbf{X}\) such that \(\alpha(x) \neq 0\) is called the support of \(\alpha\) ; then \(S_0 = \varnothing\) and \(S_{\alpha + \beta} \subset S_{\alpha} \cup S_{\beta}\) for \(\alpha, \beta\) in \(\mathbf{Z}^{\mathbf{x}}\) . Therefore the set \(\mathbf{Z}^{(\mathbf{x})}\) of mappings \(\alpha: \mathbf{X} \to \mathbf{Z}\) of finite support is a subgroup of \(\mathbf{Z}^{\mathbf{x}}\) called the free commutative group constructed on \(\mathbf{X}\) .

For all \(x \in \mathbf{X}\) , let \(\delta_x\) denote the element of \(\mathbf{Z}^{(\mathbf{x})}\) defined by

\[
\delta_ {x} (y) = \left\{ \begin{array}{l l} 1 & \text { if } y = x \\ 0 & \text { if } y \neq x. \end{array} \right.\tag{10}
\]

Also, for \(\alpha \in \mathbf{Z}^{(\mathbf{x})}\) , the integer \(|\alpha|\) , the length of \(\alpha\) is defined by the formula

\[
| \alpha | = \sum_ {x \in X} \alpha (x).\tag{11}
\]

The following relations are immediately established:

\[
\alpha = \sum_ {x \in \mathbf {X}} \alpha (x) \cdot \delta_ {x}\tag{12}
\]

\[
\left| \delta_ {x} \right| = 1, \quad \left| 0 \right| = 0\tag{13}
\]

\[
| \alpha + \beta | = | \alpha | + | \beta |\tag{14}
\]

for \(\alpha, \beta\) in \(\mathbf{Z}^{(\mathbf{X})}\) and \(x\) in \(\mathbf{X}\) .

The order relation \(\alpha \leqslant \beta\) is defined in \(\mathbf{Z}^{(\mathbf{x})}\) by \(\alpha(x) \leqslant \beta(x)\) for all \(x \in \mathbf{X}\) . The relations \(\alpha \leqslant \beta\) and \(\alpha' \leqslant \beta'\) imply \(\alpha + \alpha' \leqslant \beta + \beta'\) , \(|\alpha| \leqslant |\beta|\) and \(-\alpha \geqslant -\beta\) ; further, the relation \(\alpha \leqslant \beta\) is equivalent to \(\beta - \alpha \geqslant 0\) . The set of elements \(\alpha \geqslant 0\) in \(\mathbf{Z}^{(\mathbf{x})}\) is denoted by \(\mathbf{N}^{(\mathbf{x})}\) ; it is the set of mappings of X into N of finite support and it is a submonoid of \(\mathbf{Z}^{(\mathbf{x})}\) called the free commutative monoid constructed on X. The elements of length 1 are the minimal elements in \(\mathbf{N}^{(\mathbf{x})} = \{0\}\) and constitute the set of \(\delta_x\) ( \(x \in \mathbf{X}\) ).

The monoid \(\mathbf{N}^{(\mathbf{x})}\) and the group \(\mathbf{Z}^{(\mathbf{x})}\) enjoy the following universal property.

PROPOSITION 10. Let M be a commutative monoid (resp. group) and f a mapping of X into M. There exists one and only one homomorphism of \(\mathbf{N}^{(\mathbf{x})}\) (resp. \(\mathbf{Z}^{(\mathbf{x})}\) ) into M such that \(g(\delta_x) = f(x)\) for all \(x \in \mathbf{X}\) . If \(\mathbf{M}\) is written additively, then \(g(\alpha) = \sum_{x \in \mathbf{X}} \alpha(x) \cdot f(x)\) for all \(\alpha\) in \(\mathbf{N}^{(\mathbf{x})}\) (resp. \(\mathbf{Z}^{(\mathbf{x})}\) ).

Let \(g\) be a homomorphism of \(\mathbf{N}^{(\mathbf{X})}\) (resp. \(\mathbf{Z}^{(\mathbf{X})}\) ) into \(\mathbf{M}\) such that \(g(\delta_x) = f(x)\) for all \(x \in \mathbf{X}\) . For all \(\alpha\) in \(\mathbf{N}^{(\mathbf{X})}\) (resp. \(\mathbf{Z}^{(\mathbf{X})}\) ), it follows from (12) that

\[
g (\alpha) = \sum_ {x \in \mathbf {X}} \alpha (x). g (\delta_ {x}) = \sum_ {x \in \mathbf {X}} \alpha (x). f (x),
\]

whence the uniqueness of g.

For all \(\alpha\) in \(\mathbf{N}^{(\mathbf{x})}\) (resp. \(\mathbf{Z}^{(\mathbf{x})}\) ) we write \(g(\alpha) = \sum_{x\in \mathbf{X}}\alpha (x).f(x)\) . Then obviously \(g(0) = 0\) ; for \(\alpha, \beta\) in \(\mathbf{N}^{(\mathbf{x})}\) (resp. \(\mathbf{Z}^{(\mathbf{x})}\) ),

\[
\begin{array}{r l} g (\alpha + \beta) & = \sum_ {x \in \mathbf {X}} (\alpha (x) + \beta (x)). f (x) \\ & = \sum_ {x \in \mathbf {X}} [ \alpha (x). f (x) + \beta (x). f (x) ] \\ & = \sum_ {x \in \mathbf {X}} \alpha (x). f (x) + \sum_ {x \in \mathbf {X}} \beta (x). f (x) \\ & = g (\alpha) + g (\beta) \end{array}
\]

and hence \(g\) is a homomorphism of \(\mathbf{N}^{(\mathbf{x})}\) (resp. \(\mathbf{Z}^{(\mathbf{x})}\) ) into \(\mathbf{M}\) . Also, for \(y\) in \(\mathbf{X}\) ,

\[
g (\delta_ {y}) = \sum_ {x \in \mathbf {X}} \delta_ {y} (x) \cdot f (x);
\]

now \(\delta_y(x).f(x) = 0\) for \(x\neq y\) and \(\delta_y(y).f(y) = f(y)\) , whence \(g(\delta_y) = f(y)\) .

Let \(u: X \to Y\) be a mapping. By Proposition 10 there exists one and only one homomorphism of \(\mathbf{Z}^{(X)}\) into \(\mathbf{Z}^{(Y)}\) which maps \(\delta_x\) to \(\delta_{u(x)}\) for all \(x \in X\) . It is denoted by \(\mathbf{Z}^{(w)}\) ; it is immediately seen that it maps \(\alpha \in \mathbf{Z}^{(X)}\) to the element \(\beta \in \mathbf{Z}^{(Y)}\) defined by

\[
\beta (y) = \sum_ {x \in u ^ {- 1} (y)} \alpha (x).\tag{15}
\]

As in the case of magmas (no. 1), the formula \(\mathbf{Z}^{(v\circ u)} = \mathbf{Z}^{(v)}\circ \mathbf{Z}^{(y)}\) is established for every mapping \(v\colon \mathbf{Y}\to \mathbf{Z}\) ; it is also shown that \(\mathbf{Z}^{(u)}\) is injective (resp. surjective, bijective) if \(u\) is.

Let \(S\) be a subset of \(X\) ; if \(i\) is the injection of \(S\) into \(X\) , the mapping \(f = \mathbf{Z}^{(i)}\) is an isomorphism of \(\mathbf{Z}^{(S)}\) onto the subgroup \(H\) of \(\mathbf{Z}^{(X)}\) generated by the elements \(\delta_s\) for \(s \in S\) . By (15),

\[
(f (\alpha)) (x) = \left\{ \begin{array}{l l} \alpha (x) & \text { if } x \in S \\ 0 & \text { if } x \in X - S \end{array} \right.
\]

and therefore H is the set of elements of \(\mathbf{Z}^{(\mathbf{x})}\) of support contained in S. Henceforth \(\mathbf{Z}^{(\mathbf{s})}\) will be identified with H by means of f.

Formula (15) shows that the restriction of \(\mathbf{Z}^{(u)}\) to \(\mathbf{N}^{(\mathbf{X})}\) induces a homomorphism \(\mathbf{N}^{(u)}\) of \(\mathbf{N}^{(\mathbf{X})}\) into \(\mathbf{N}^{(\mathbf{Y})}\) . Then \(\mathbf{N}^{(v\circ u)} = \mathbf{N}^{(v)}\circ \mathbf{N}^{(u)}\) for every mapping \(v\colon \mathbf{Y}\to \mathbf{Z}\) ; further, \(\mathbf{N}^{(u)}\) is injective (resp. surjective, bijective) if \(u\) is. If \(S\) is a subset of \(\mathbf{X}\) ,

\[
\mathbf {N} ^ {(\mathbf {S})} = \mathbf {Z} ^ {(\mathbf {S})} \cap \mathbf {N} ^ {(\mathbf {X})}.
\]

Remark. Let M be the multiplicative monoid of strictly positive integers and let P be the set of prime numbers (§ 4, no. 10, Definition 15). By Proposition 10 there exists a homomorphism u of \(\mathbf{N}^{(\mathfrak{P})}\) into M characterized by \(u(\delta_{p}) = p\) for every prime number p.~Then \(u(\alpha) = \prod_{p \in \mathfrak{P}} p^{\alpha(p)}\) for \(\alpha\) in \(\mathbf{N}^{(\mathfrak{P})}\) and Theorem 7 of § 4, no. 10 shows that u is an isomorphism of \(\mathbf{N}^{(\mathfrak{P})}\) onto H.

\subsection{EXPONENTIAL NOTATION}

Let M be a monoid, written multiplicatively, and \(\mathbf{u} = (u_{x})_{x \in \mathbf{X}}\) a family of elements of M, commuting in pairs. Let \(\alpha\) be in \(\mathbf{N}^{(\mathbf{X})}\) ; the elements \(u_{x}^{\alpha(x)}\) and \(u_{y}^{\alpha(y)}\) of M commute for x, y in X and there exists a finite subset S of X such that \(u_{x}^{\alpha(x)} = 1\) for x in X - S. We may therefore write:

\[
\mathbf {u} ^ {\alpha} = \prod_ {x \in \mathrm{X}} u _ {x} ^ {\alpha (x)}.\tag{16}
\]

Let \(\mathbf{M}'\) be the submonoid of \(\mathbf{M}\) generated by the family \((u_x)_{x\in \mathbf{X}}\) ; it is commutative (§ 1, no. 5, Corollary 2 to Proposition 4). There thus exists (no. 7, Proposition 10) a unique homomorphism \(f\) of \(\mathbf{N}^{(\mathbf{X})}\) into \(\mathbf{M}'\) such that \(f(\delta_x) = u_x\) for all \(x\in \mathbf{X}\) and \(f(\alpha) = \mathbf{u}^{\alpha}\) for all \(\alpha\) in \(\mathbf{N}^{(\mathbf{X})}\) . We deduce the following formulae

\[
\mathbf {u} ^ {\alpha + \beta} = \mathbf {u} ^ {\alpha}. \mathbf {u} ^ {\beta}\tag{17}
\]

\[
\mathbf {u} ^ {0} = 1\tag{18}
\]

\[
\mathbf {u} ^ {\delta_ {x}} = u _ {x}\tag{19}
\]

for \(\alpha, \beta\) in \(\mathbf{N}^{(\mathbf{x})}\) and \(x\) in \(\mathbf{X}\) .

Let \(\mathbf{v} = (v_x)_{x \in \mathbf{X}}\) be another family of elements of \(\mathbf{M}\) ; suppose that \(v_x v_y = v_y v_x\) and \(u_x v_y = v_y u_x\) for \(x, y\) in \(\mathbf{X}\) . Then there exists (§ 1, no. 5, Corollary 2 to Proposition 4) a commutative submonoid \(\mathbf{L}\) of \(\mathbf{M}\) such that \(u_x \in \mathbf{L}\) and \(v_x \in \mathbf{L}\) for all \(x \in \mathbf{X}\) . The mapping \(\alpha \mapsto u^\alpha . v^\alpha\) of \(\mathbf{N}^{(\mathbf{X})}\) into \(\mathbf{L}\) is then a homomorphism (§ 1, no. 5, Proposition 5) mapping \(\delta_x\) to \(u_x . v_x\) . Thus we have the formula

\[
\mathbf {u} ^ {\alpha}. \mathbf {v} ^ {\alpha} = (\mathbf {u}. \mathbf {v}) ^ {\alpha},\tag{20}
\]

where \(\mathbf{u}.\mathbf{v}\) is the family \((u_x.v_x)_{x\in \mathbf{X}}\)

When M is commutative, \(u^{\alpha}\) can be defined for every family u of elements of M and formulae (15) to (20) hold without restriction.

\subsection{RELATIONS BETWEEN THE VARIOUS FREE OBJECTS}

As the free monoid \(\mathrm{Mo}(\mathbf{X})\) is a magma, Proposition 1, of no. 1 shows the existence of a homomorphism \(\lambda: \mathrm{M}(\mathbf{X}) \to \mathrm{Mo}(\mathbf{X})\) whose restriction to \(\mathbf{X}\) is the identity. Similarly, as the free group \(\mathrm{F}(\mathbf{X})\) is a monoid the identity mapping of \(\mathbf{X}\) extends to a homomorphism \(\mu: \mathrm{Mo}(\mathbf{X}) \to \mathrm{F}(\mathbf{X})\) (no. 2, Proposition 3). By no. 4, Proposition 6 and no. 5, Proposition 7, \(\mu\) is injective. Similarly Proposition 10 of no. 7 and Proposition 8 of no. 5 show the existence of homomorphisms \(\nu: \mathrm{Mo}(\mathbf{X}) \to \mathbf{N}^{(\mathbf{X})}\) and \(\pi: \mathrm{F}(\mathbf{X}) \to \mathbf{Z}^{(\mathbf{X})}\) characterized by \(\nu(x) = \delta_x\) and \(\pi(x) = \delta_x\) for all \(x \in \mathbf{X}\) . If \(\iota\) is the injection of \(\mathbf{N}^{(\mathbf{X})}\) into \(\mathbf{Z}^{(\mathbf{X})}\) , the two homomorphisms \(\iota \circ \nu\) and \(\pi \circ \mu\) of \(\mathrm{Mo}(\mathbf{X})\) into \(\mathbf{Z}^{(\mathbf{X})}\) coincide on \(\mathbf{X}\) , whence \(\iota \circ \nu = \pi \circ \mu\) . The situation may be summarized by the following commutative diagram:

\[
\begin{array}{c} \mathrm{M} (\mathrm{X}) \xrightarrow {\lambda} \mathrm{Mo} (\mathrm{X}) \xrightarrow {\nu} \mathbf {N} ^ {(\mathrm{X})} \\ \mu \Biggl \downarrow \\ \mathrm{F} (\mathrm{X}) \xrightarrow {\pi} \mathbf {Z} ^ {(\mathrm{X})}. \end{array}
\]

The homomorphisms \(\lambda\) , \(\mu\) , \(\nu\) and \(\pi\) will be called canonical.

Let \(w\) be in \(\mathbf{M}(\mathbf{X})\) ; it is immediately shown by induction on \(l(w)\) that the length of the word \(\lambda(w)\) is equal to that of \(w\) . Moreover

\[
\nu (x _ {1} \dots x _ {n}) = \sum_ {i = 1} ^ {n} \delta_ {x _ {i}}\tag{21}
\]

for \(x_1, \ldots, x_n\) in X, whence \(|\nu(x_1 \ldots x_n)| = n\) by (13) and (14). In other words,

\[
\left| \nu (u) \right| = l (u) \quad (u \in \operatorname{Mo} (\mathbf {X})).\tag{22}
\]

PROPOSITION 11. The canonical homomorphism \(\nu\) of \(\operatorname{Mo}(\mathbf{X})\) into \(\mathbf{N}^{(\mathbf{X})}\) is surjective. Let \(w = x_1 \ldots x_n\) and \(w' = x_1' \ldots x_m'\) be two elements of \(\operatorname{Mo}(\mathbf{X})\) ; in order that \(\nu(w) = \nu(w')\) , it is necessary and sufficient that \(m = n\) and that there exist a permutation \(\sigma \in \mathfrak{S}_n\) with \(x_i' = x_{\sigma(i)}\) for \(1 \leqslant i \leqslant n\) .

The image of \(\nu\) is a submonoid I of \(\mathbf{N}^{(\mathbf{x})}\) containing the elements \(\delta_x\) (for \(x \in \mathbf{X}\) ). Formula (12) (no. 7) shows that \(\mathbf{N}^{(\mathbf{x})}\) is generated by the family \((\delta_x)_{x \in \mathbf{X}}\) , where \(\mathbf{I} = \mathbf{N}^{(\mathbf{x})}\) . Therefore \(\nu\) is surjective.

If \(m = n\) and \(x_{i}^{\prime} = x_{\sigma (i)}\) for \(\mathrm{l}\leqslant i\leqslant n\) , then

\[
\nu (w ^ {\prime}) = \sum_ {i = 1} ^ {n} \delta_ {x _ {i} ^ {\prime}} = \sum_ {i = 1} ^ {n} \delta_ {x _ {\sigma (i)}} = \sum_ {i = 1} ^ {n} \delta_ {x _ {i}} = \nu (w)
\]

by formula (21) and the commutativity theorem (§ 1, no. 5, Theorem 2).

Conversely, suppose that \(\nu(w)\) and \(\nu(w')\) are equal to the same element \(\alpha\) of \(\mathbf{N}^{(\mathbf{x})}\) ; by formula (22), \(n = |\alpha| = m\) . For all \(x \in \mathbf{X}\) , let \(\mathrm{I}_x\) (resp. \(\mathrm{I}'_x\) ) be the set of integers \(i\) such that \(1 \leqslant i \leqslant n\) and \(x_i = x\) (resp. \(x_i' = x\) ). Hence \((\mathrm{I}_x)_{x \in \mathbf{X}}\) and \((\mathrm{I}'_x)_{x \in \mathbf{X}}\) are partitions of the interval \([1, n]\) of \(\mathbf{N}\) ; further, the formula \(\alpha = \sum_{i=1}^{n} \delta_{x_i}\) shows that \(\alpha(x)\) is the cardinal of \(I_x\) ; similarly the formula \(\alpha = \sum_{i=1}^{\infty} \delta_{x'_i}\) shows that \(\alpha(x)\) is the cardinal of \(\mathbf{I}_x'\) . There thus exists a permutation \(\sigma\) of \([1, n]\) such that \(\sigma(\mathbf{I}_x') = \mathbf{I}_x\) for all \(x \in \mathbf{X}\) , that is \(x'_i = x_{\sigma(i)}\) for \(i = 1, \ldots, n\) .

Remark. Let S be a subset of X. Recall that we have identified M(S) with a submagma of M(X), Mo(S) with a submonoid of Mo(X) and \(\mathbf{N}^{(S)}\) with a submonoid of \(\mathbf{N}^{(X)}\) . Then

\[
\mathbf {M} (\mathrm{S}) = \lambda^ {- 1} (\mathbf {M o} (\mathrm{S})).\tag{23}
\]

Clearly \(\lambda (\mathbf{M}(\mathbf{S}))\subset \mathbf{Mo}(\mathbf{S})\) . Let \(w\in \lambda^{-1}(\mathbf{Mo}(\mathbf{S}))\) ; we show by induction on \(l(w)\) that \(w\in \mathbf{M}(\mathbf{S})\) . It is obvious if \(l(w) = 1\) . If \(l(w) > 1\) , we may write \(w = w_{1}w_{2}\) with \(w_{1},w_{2}\in \mathbf{M}(\mathbf{X}), l(w_{1}) <   l(w),l(w_{2}) <   l(w)\) . Then \(\lambda (w_1)\lambda (w_2)\in \mathbf{Mo}(\mathbf{S})\) hence \(\lambda (w_1)\in \mathbf{Mo}(\mathbf{S})\) and \(\lambda (w_2)\in \mathbf{Mo}(\mathbf{S})\) , whence \(w_{1}\in \mathbf{M}(\mathbf{S})\) and \(w_{2}\in \mathbf{M}(\mathbf{S})\) by the induction hypothesis and finally \(w\in \mathbf{M}(\mathbf{S})\)

Also

\[
\mathbf {M o} (\mathrm{S}) = \nu^ {- 1} (\mathbf {N} ^ {(\mathrm{S})}).\tag{24}
\]

This follows immediately from formula (21).

Further, \(\mathbf{N}^{(\mathbf{S})}\) is the set of elements of \(\mathbf{N}^{(\mathbf{X})}\) whose support is contained in S; if \((\mathrm{S}_i)_{i\in I}\) is a family of subsets of X of intersection S, then \(\mathbf{N}^{(\mathbf{S})} = \bigcap_{i\in I}\mathbf{N}^{(\mathbf{S}_i)}\) and formulae (23) and (24) imply

\[
\mathbf {M} (\mathrm{S}) = \bigcap_ {i \in \mathrm{I}} \mathbf {M} (\mathrm{S} _ {i}), \quad \mathbf {M o} (\mathrm{S}) = \bigcap_ {i \in \mathrm{I}} \mathbf {M o} (\mathrm{S} _ {i}).\tag{25}
\]

\section{RINGS}

\subsection{RINGS}

DEFINITION 1. A ring is a set A with two laws of composition called respectively addition and multiplication, satisfying the following axioms:

(AN I) Under addition A is a commutative group.

(AN II) Multiplication is associative and possesses an identity element.

(AN III) Multiplication is distributive with respect to addition.

The ring A is said to be commutative if its multiplication is commutative.

In what follows, \((x,y)\mapsto x+y\) denotes addition and \((x,y)\mapsto xy\) multiplication; 0 denotes the identity element for addition and 1 that for multiplication. Finally, -x denotes the negative of x under addition. The axioms of a ring are therefore expressed by the following identities:

\[
x + (y + z) = (x + y) + z \quad (\text { associativity   of   addition })\tag{1}
\]

\[
x + y = y + x \quad (\text { commutativity   of   addition })\tag{2}
\]

\[
0 + x = x + 0 = x \quad (\text { zero })\tag{3}
\]

\[
x + (- x) = (- x) + x = 0 \quad (\text { negative })\tag{4}
\]

\[
x (y z) = (x y) z \quad (\text { associativity   of   multiplication })\tag{5}
\]

\[
x. 1 = 1. x = x \quad (\text { unit   element })\tag{6}
\]

\[
\left. \begin{array}{l} (x + y). z = x z + y z \\ x. (y + z) = x y + x z \end{array} \right\} (\text {distributivity})\tag{7, 8}
\]

Finally the ring A is commutative if \(xy = yx\) for \(x, y\) in A.

With addition alone A is a commutative group called the additive group of A. For all \(x \in A\) , we define the left homothety \(\gamma_{x}\) and right homothety \(\delta_{x}\) by \(\gamma_{x}(y) = xy\) , \(\delta_{x}(y) = yx\) . By formulae (7) and (8), \(\gamma_{x}\) and \(\delta_{x}\) are endomorphisms of the additive group of A and thus map zero to zero and negative to negative. Therefore

\[
x. 0 = 0. x = 0\tag{9}
\]

\[
x. (- y) = (- x). y = - x y;\tag{10}
\]

it follows that \((-x)(-y) = -((-x).y) = -(-xy)\) , whence

\[
(- x) (- y) = x y.\tag{11}
\]

Formulae (10) and (11) constitute the sign rule. It follows that

\[
- x = (- 1) x = x (- 1)
\]

and \((-1)(-1) = 1\) .

From (11) it follows by induction on \(n\) that

\[
(- x) ^ {n} = \left\{ \begin{array}{l l} x ^ {n} & \text { if   } n \text {   is   even } \\ - x ^ {n} & \text { if   } n \text {   is   odd. } \end{array} \right.\tag{12}
\]

When we speak of cancellable elements, invertible elements, permutable elements, central elements, centralizer or centre of a ring A, all these notions will refer to the multiplication on A. If \(x, y \in A\) and y is invertible, the element \(xy^{-1}\) of A is also denoted by x/y when A is commutative. The set of invertible elements of A is stable under multiplication. Under the law induced by multiplication it is a group called the multiplicative group of A, sometimes denoted by A*.

Let \(x, y\) be in A. \(x\) is said to be a left (resp. right) multiple of \(y\) if there exists \(y' \in \mathbf{A}\) such that \(x = y'y\) (resp. \(x = yy'\) ); it is also said that \(y\) is a right (resp.

left) divisor of x. When A is commutative, there is no need to distinguish between ``left'' and ``right''.

In conformity with the above terminology, every element \(y \in A\) would be considered as a right and left divisor of 0; but, by an abuse of language, in general the term ``right (resp. left) divisor of 0'' is reserved for elements y such that there exists \(x \neq 0\) in A satisfying the relation xy = 0 (resp. yx = 0). In other words, the right (resp. left) divisors of zero are the right (resp. left) non-cancellable elements.

Let \(x \in A\) . \(x\) is called nilpotent if there exists an integer \(n > 0\) with \(x^n = 0\) . The element \(1 - x\) is then invertible, with inverse equal to

\[
1 + x + x ^ {2} + \dots + x ^ {n - 1}.
\]

As A is a commutative group under addition, the element \(nx\) for \(n \in \mathbf{Z}\) and \(x \in \mathbf{A}\) has been defined (§ 2, no. 8). As \(\gamma_x\) and \(\delta_x\) are endomorphisms of the additive group \(\mathbf{A}\) , \(\gamma_x(ny) = n\gamma_x(y)\) and \(\delta_y(nx) = n\delta_y(x)\) , whence

\[
x. (n y) = (n x). y = n. (x y).
\]

In particular, \(nx = (n.1)x\) .

A set A with addition and multiplication satisfying the axioms of a ring with the exception of that assuring the existence of the identity element under multiplication, is called a pseudo-ring.

\subsection{CONSEQUENCES OF DISTRIBUTIVITY}

Distributivity of multiplication with respect to addition allows us to apply Proposition 1 of §3, no. 4, which gives

\[
\prod_ {i = 1} ^ {n} \left(\sum_ {\lambda \in L _ {i}} x _ {i, \lambda}\right) = \sum_ {\alpha_ {1}, \dots , \alpha_ {n}} \prod_ {i = 1} ^ {n} x _ {i, \alpha_ {i}}\tag{13}
\]

where the sum extends over all sequences \((\alpha_{1},\ldots,\alpha_{n})\) belonging to \(L_{1}\times\cdots\times L_{n}\) and for \(i=1,\ldots,n\) the family \((x_{i,\lambda})_{\lambda\in L^{i}}\) of elements of the ring A is of finite support.

PROPOSITION 1. Let A be a commutative ring and \((x_{\lambda})_{\lambda \in L}\) a finite family of elements of A. For every family of positive integers \(\beta = (\beta_{\lambda})_{\lambda \in L}\) , let \(|\beta| = \sum_{\lambda \in L} \beta_{\lambda}\) . Then

\[
\left(\sum_ {\lambda \in L} x _ {\lambda}\right) ^ {n} = \sum_ {| \beta | = n} \frac {n !}{\prod_ {\lambda \in L} \beta_ {\lambda} !} \prod_ {\lambda \in L} x _ {\lambda} ^ {\beta_ {\lambda}}.\tag{14}
\]

We apply formula (13) with \(\mathbf{L}_i = \mathbf{L}\) and \(x_{i,\lambda} = x_{\lambda}\) for \(1 \leqslant i \leqslant n\) . Then

\[
\left(\sum_ {\lambda \in L} x _ {\lambda}\right) ^ {n} = \sum_ {\alpha_ {1}, \dots , \alpha_ {n}} x _ {\alpha_ {1}} \dots x _ {\alpha_ {n}},\tag{15}
\]

the sum extending over all sequences \(\alpha = (\alpha_{1},\ldots ,\alpha_{n})\in \mathbf{L}^{n}\)

Let \(\alpha\) be in \(\mathbf{L}^n\) ; for all \(\lambda \in \mathbf{L}\) let \(\mathrm{U}_{\lambda}^{\alpha}\) denote the set of integers \(i\) such that \(1 \leqslant i \leqslant n\) and \(\alpha_i = \lambda\) and let \(\Phi(\alpha) = (\mathrm{U}_{\lambda}^{\alpha})_{\lambda \in \mathbf{L}}\) . It is immediate that \(\Phi\) is a bijection of \(\mathbf{L}^n\) onto the set of partitions of \(\{1, 2, \ldots, n\}\) indexed by \(\mathbf{L}\) . For all \(\beta \in \mathbf{N}^{\mathbf{L}}\) such that \(|\beta| = n\) , let \(\mathbf{L}_{\beta}^n\) denote the set of \(\alpha \in \mathbf{L}^n\) such that \(\operatorname{Card} \mathrm{U}_{\beta}^{\alpha} = \beta_{\lambda}\) for all \(\lambda \in \mathbf{L}\) . It follows that the family \((\mathbf{L}_{\beta}^n)_{\beta | = n}\) is a partition of \(\mathbf{L}^n\) and that

\[
\operatorname{Card} \mathrm{L} _ {\beta} ^ {n} = \frac {n !}{\prod_ {\lambda \in \mathrm{L}} \beta_ {\lambda} !}
\]

(§ 5, no. 5).

Finally, for \(\alpha \in \mathbf{L}_{\beta}^{n}\) ,

\[
x _ {\alpha_ {1}} \dots x _ {\alpha_ {n}} = \prod_ {\lambda \in L} \prod_ {i \in U _ {\lambda} ^ {\alpha}} x _ {\alpha_ {i}} = \prod_ {\lambda \in L} \prod_ {i \in U _ {\lambda} ^ {\alpha}} x _ {\lambda} = \prod_ {\lambda \in L} x _ {\lambda},
\]

whence

\[
\begin{array}{r l} \sum_ {(\alpha_ {1}, \dots , \alpha_ {n})} x _ {\alpha_ {1}} \dots x _ {\alpha_ {n}} & = \sum_ {| \beta | = n} \sum_ {\alpha \in L _ {\beta} ^ {n}} x _ {\alpha_ {1}} \dots x _ {\alpha_ {n}} \\ & = \sum_ {| \beta | = n} \sum_ {\alpha \in L _ {\beta} ^ {n}} \prod_ {\lambda \in L _ {\lambda}} x _ {\lambda} ^ {\beta_ {\lambda}} \\ & = \sum_ {| \beta | = n} \frac {n !}{\prod_ {\lambda \in L} \beta_ {\lambda} !} x _ {\lambda} ^ {\beta_ {\lambda}} \end{array}
\]

and formula (14) thus follows from (15).

COROLLARY 1 (binomial formula). Let \(x\) and \(y\) be two elements of a commutative ring A. Then:

\[
(x + y) ^ {n} = \sum_ {p = 0} ^ {n} \binom {n} {p} x ^ {p} y ^ {n - p}.
\]

Formula (14) applied to \(\mathbf{L} = \{1,2\}\) , \(x_{1} = x\) and \(x_{2} = y\) gives

\[
(x + y) ^ {n} = \sum_ {p + q = n} \frac {n !}{p ! q !} x ^ {p} y ^ {q},
\]

the sum extending over ordered pairs of positive integers \(p, q\) with \(p + q = n\) . The binomial formula follows immediately from this (Set Theory, III, § 5, no. 8).

COROLLARY 2. Let A be a commutative ring, X a set, \(\mathbf{u} = (u_x)_{x \in \mathbb{X}}\) and \(\mathbf{v} = (v_x)_{x \in \mathbb{X}}\) two families of elements of A. Let \(\mathbf{u} + \mathbf{v}\) denote the family \((u_x + v_x)_{x \in \mathbb{X}}\) . For all \(\lambda \in \mathbf{N}^{(\mathbb{X})}\) we write \(\lambda! = \prod_{x \in \mathbb{X}} \lambda(x)!\) . Then for all \(\alpha \in \mathbf{N}^{(\mathbb{X})}\) , in the notation of § 7, no. 8,

\[
(\mathbf {u} + \mathbf {v}) ^ {\alpha} = \sum_ {\beta + \gamma = \alpha} \frac {\alpha !}{\beta ! \gamma !} \mathbf {u} ^ {\beta} \mathbf {v} ^ {\gamma}.
\]

For \(x\in \mathbf{X}\)

\[
(u _ {x} + v _ {x}) ^ {\alpha (x)} = \sum_ {m + n = \alpha (x)} \frac {\alpha (x) !}{m ! n !} u _ {x} ^ {m} v _ {x} ^ {n}
\]

by Corollary 1. Taking the product of these equations for \(x \in \mathbf{X}\) and using (13), we obtain the corollary.

PROPOSITION 2. Let A be a ring, \(x_{1}, \ldots, x_{n}\) elements of A and I = \{1, 2, \ldots, n\}. For H ⊂ I, we write \(x_{\mathrm{H}} = \sum_{i \in \mathrm{H}} x_{i}\) . Then

\[
(- 1) ^ {n} \sum_ {\sigma \in \mathfrak {S} _ {n}} x _ {\sigma (1)} \dots x _ {\sigma (n)} = \sum_ {H \subset I} (- 1) ^ {\operatorname{Card} H} (x _ {H}) ^ {n}.\tag{16}
\]

In particular, if A is commutative,

\[
(- 1) ^ {n} n! x _ {1} x _ {2} \dots x _ {n} = \sum_ {\mathrm{H} \in \mathrm{I}} (- 1) ^ {\text { Card   H }} (x _ {\mathrm{H}}) ^ {n}.
\]

Let C be the set of mappings of I into \(\{0,1\}\) . If each \(H \subset I\) is mapped to its characteristic function, a bijection is obtained of \(\mathfrak{P}(I)\) onto C. The right hand side of (16) is thus equal to:

\[
\begin{array}{r l} \sum_ {a \in \mathbb {C}} (- 1) ^ {a (1) + \dots + a (n)} & \left(\sum_ {i \in \mathbb {I}} a (i) x _ {i}\right) ^ {n} \\ & = \sum_ {a \in \mathbb {C}} (- 1) ^ {a (1) + \dots + a (n)} \sum_ {(i _ {1}, \ldots , i _ {n}) \in \mathbb {I} ^ {n}} a (i _ {1}) \ldots a (i _ {n}) x _ {i _ {1}} \ldots x _ {i _ {n}} \\ & = \sum_ {(i _ {1}, \ldots , i _ {n}) \in \mathbb {I} ^ {n}} c _ {i _ {1} \ldots i _ {n}} x _ {i _ {1}} \ldots x _ {i _ {n}} \end{array}
\]

where

\[
c _ {i _ {1} \dots i _ {n}} = \sum_ {a \in \mathbb {C}} (- 1) ^ {a (1) + \dots + a (n)} a (i _ {1}) \dots a (i _ {n}).
\]

\begin{enumerate}
\def\labelenumi{(\arabic{enumi})}
\tightlist
\item
  Suppose that \((i_{1},\ldots,i_{n})\) is not a permutation of I. There exists a \(j\in I\) distinct from \(i_{1},\ldots,i_{n}\) . Let \(C'\) be the set of \(a\in C\) such that \(a(j)=0\) . For all \(a\in C'\) , let \(a^{*}\) be the sum of a and the characteristic function of \(\{j\}\) . Then \(a^{*}(1)+\cdots+a^{*}(n)=a(1)+\cdots+a(n)+1\) and hence
\end{enumerate}

\[
\begin{array}{l} c _ {i _ {1} \dots i _ {n}} = \sum_ {a \in C ^ {\prime}} (- 1) ^ {a (1) + \dots + a (n)} a (i _ {1}) \dots a (i _ {n}) + (- 1) ^ {a ^ {*} (1) + \dots + a ^ {*} (n)} a ^ {*} (i _ {1}) \dots a ^ {*} (i _ {n}) \\ = \sum_ {a \in C ^ {\prime}} ((- 1) ^ {a (1) + \dots + (n)} + (- 1) ^ {a (1) + \dots + a (n) + 1}) a (i _ {1}) \dots a (i _ {n}) = 0. \end{array}
\]

\begin{enumerate}
\def\labelenumi{(\arabic{enumi})}
\setcounter{enumi}{1}
\tightlist
\item
  Suppose that there exists \(\sigma \in \mathfrak{S}_n\) such that \(i_1 = \sigma(1), \ldots, i_n = \sigma(n)\) . Then \(a(i_1) \ldots a(i_n) = 0\) unless \(a\) only takes the value 1. Thus \(c_{i_1 \ldots i_n} = (-1)^n\) .
\end{enumerate}

\subsection{EXAMPLES OF RINGS}

I. Zero ring. Let A be a ring. For 0 = 1 in A, it is necessary and sufficient that A consist of a single element. The condition is obviously sufficient. On the other hand, if 0 = 1, then, for all \(x \in A\) , x = x.1 = x.0 = 0. Such a ring is called a zero ring.

\begin{enumerate}
\def\labelenumi{\Roman{enumi}.}
\setcounter{enumi}{1}
\tightlist
\item
  Ring of rational integers. With the addition defined in §2, no. 5, and the multiplication defined in §2, no. 6, Z is a commutative ring. The notation 0, 1, -x is in accordance with the notation introduced earlier.
\end{enumerate}

*III. Ring of real-valued functions. Let I be an interval in the set R of real numbers and let A be the set of continuous functions defined on I with real values. The sum \(f + g\) and product f.g of two functions f and g are defined by

\[
(f + g) (t) = f (t) + g (t), \quad (f g) (t) = f (t) g (t) \quad (t \in \mathbf {I}).
\]

A commutative ring is obtained whose unit element is the constant 1.*

*IV. Convolution pseudo-ring. Let E be the set of real-valued continuous functions on R, which are zero outside a bounded interval. The sum of two functions is defined as in III, but the product is now defined by

\[
(f g) (t) = \int_ {- \infty} ^ {\infty} f (s) g (t - s) d s
\]

(``convolution product''). Thus a commutative pseudo-ring is obtained which is not a ring (cf.~Integration, VIII, §4).*

V. Opposite ring of a ring A. Let A be a ring. The set A with the same addition as A and the multiplication \((x,y)\mapsto yx\) is often denoted by \(A^{0}\) . It is a ring (called the opposite ring of A) with the same zero and same unit as A and which coincides with A if and only if A is commutative.

\begin{enumerate}
\def\labelenumi{\Roman{enumi}.}
\setcounter{enumi}{5}
\tightlist
\item
  Endomorphism ring of a commutative group. Let G be a commutative group written additively. Let E denote the set of endomorphisms of G. Given f and g in E, the mappings \(f + g\) and fg of G into G are defined by
\end{enumerate}

\[
(f + g) (x) = f (x) + g (x), \quad (f g) (x) = f (g (x)) \quad (x \in \mathbf {G}).
\]

By § 1, no. 5, Proposition 5, \(f + g\) is an endomorphism of G and so obviously also is \(fg = f \circ g\) . By § 4, no. 8, E is a (commutative) group under addition. Multiplication is obviously associative and has identity element \(\operatorname{Id}_{\mathbf{G}}\) . Also for \(f, g\) and \(h\) in E, we write \(\phi = f \cdot (g + h)\) ; for all \(x \in \mathbf{G}\) ,

\[
\phi (x) = f ((g + h) (x)) = f (g (x) + h (x)) = f (g (x)) + f (h (x))
\]

for \(f\) is an endomorphism of \(\mathbf{G}\) ; hence \(\phi = fg + fh\) and clearly

\[
(g + h) f = g f + h f.
\]

Therefore E is a (not in general commutative) ring called the endomorphism ring of G.

\begin{enumerate}
\def\labelenumi{\Roman{enumi}.}
\setcounter{enumi}{6}
\tightlist
\item
  Pseudo-ring of zero square. A pseudo-ring A is said to be of zero square if xy = 0 for all \(x, y \in A\) . Let G be a commutative group. If the set G is given the addition of the group G and multiplication \((x, y) \mapsto 0\) , a pseudo-ring of zero square is obtained. It is only a ring if \(G = \{0\}\) , in which case it is the zero ring.
\end{enumerate}

\subsection{RING HOMOMORPHISMS}

DEFINITION 2. Let A and B be two rings. A morphism, or homomorphism, of A into B is any mapping f of A into B satisfying the relations:

\[
f (x + y) = f (x) + f (y), \quad f (x y) = f (x). f (y), \quad f (1) = 1,\tag{17}
\]

for all \(x, y\) in A.

The composition of two ring homomorphisms is a ring homomorphism. Let A and B be two rings and f a mapping of A into B; for f to be an isomorphism, it is necessary and sufficient that it be a bijective homomorphism; in that case, \(\overline{f}^{1}\) is a homomorphism of B into A. A homomorphism of a ring A into itself is called an endomorphism of A.

Let \(f: A \to B\) be a ring homomorphism. The mapping \(f\) is a homomorphism of the additive group of \(A\) into the additive group of \(B\) ; in particular, \(f(0) = 0\) and \(f(-x) = -f(x)\) for all \(x \in A\) . The image under \(f\) of an invertible element of \(A\) is an invertible element of \(B\) and \(f\) induces a homomorphism of the multiplicative group of \(A\) into the multiplicative group of \(B\) .

Examples. (1) Let A be a ring. It is immediately seen that the mapping \(n \mapsto n\) . 1 of Z into A is the unique homomorphism of Z into A. In particular, the identity mapping of Z is the unique endomorphism of the ring Z.

In particular, take A to be the endomorphism ring of the additive group Z (no. 3, Example VI). The mapping \(n \mapsto n\) . 1 of Z into A is an isomorphism of Z onto A by the very construction of multiplication in Z (§ 2, no. 6).

\begin{enumerate}
\def\labelenumi{(\arabic{enumi})}
\setcounter{enumi}{1}
\tightlist
\item
  Let \(a\) be an invertible element of a ring \(A\) . The mapping \(x \mapsto axa^{-1}\) is an endomorphism of \(A\) for
\end{enumerate}

\[
\begin{array}{r l} a (x + y) a ^ {- 1} & = a x a ^ {- 1} + a y a ^ {- 1}, \\ a (x y) a ^ {- 1} & = (a x a ^ {- 1}) (a y a ^ {- 1}). \end{array}
\]

It is bijective, for the relation \(x' = axa^{-1}\) is equivalent to \(x = a^{-1}x'a\) . It is therefore an automorphism of the ring A, called the inner automorphism associated with a.

\subsection{SUBRINGS}

DEFINITION 3. Let A be a ring. A subring of A is any subset B of A which is a subgroup of A under addition, which is stable under multiplication and which contains the unit of A.

The above conditions may be written as follows

\[
0 \in \mathbf {B}, \quad \mathbf {B} + \mathbf {B} \subset \mathbf {B}, \quad - \mathbf {B} \subset \mathbf {B}, \quad \mathbf {B}. \mathbf {B} \subset \mathbf {B}, \quad 1 \in \mathbf {B}.
\]

If B is a subring of A, it is given the addition and multiplication induced by those on A, which make it into a ring. The canonical injection of B into A is a ring homomorphism.

Examples. (1) Every subgroup of the additive group \(\mathbf{Z}\) which contains 1 is equal to \(\mathbf{Z}\) . Thus \(\mathbf{Z}\) is the only subring of \(\mathbf{Z}\) .

\begin{enumerate}
\def\labelenumi{(\arabic{enumi})}
\setcounter{enumi}{1}
\item
  Let A be a ring and \((A_{t})_{t\in I}\) a family of subrings of A; it is immediate that \(\bigcap_{i\in I}A_{i}\) is a subring of A. In particular, the intersection of the subrings of A containing a subset X of A is a subring called the subring of A generated by X.
\item
  Let X be a subset of a ring A. The centralizer of X in A is a subring of A. In particular, the centre of A is a subring of A.
\item
  Let G be a commutative group with operators; let \(\Omega\) denote the set of operators and \(\alpha \mapsto f_{\alpha}\) the action of \(\Omega\) on G. Let E be the endomorphism ring of the group without operators G and F the set of endomorphisms of the group with operators G. By definition, F consists of the endomorphisms \(\phi\) of G such that \(\phi \cdot f_{\alpha} = f_{\alpha} \cdot \phi\) for all \(\alpha \in \Omega\) . Therefore F is a subring of the ring E. F is called the endomorphism ring of the group with operators G (cf.~II, § 1, no. 2). Let \(F_{1}\) be the subring of E generated by the \(f_{\alpha}\) . Then F is the centralizer of \(F_{1}\) in E.
\end{enumerate}

\subsection{IDEALS}

DEFINITION 4. Let A be a ring. A subset \(\mathfrak{a}\) of A is called a left (resp. right) ideal if it is a subgroup of the additive group of A and the relations \(a \in \mathbf{A}\) , \(x \in \mathfrak{a}\) imply \(ax \in \mathfrak{a}\) (resp. \(xa \in \mathfrak{a}\) ). \(\mathfrak{a}\) is called a two-sided ideal of A if it is both a left ideal and a right ideal of A.

The definition of a left ideal may be expressed by the relations

\[
0 \in a, \quad a + a \subset a, \quad A. a \subset a
\]

the relation \(-a \subset a\) following from the formula \((-1).x = -x\) and \(A.a \subset a\) . For all \(x \in A\) , let \(\gamma_x\) be the mapping \(a \mapsto xa\) of \(A\) into \(A\) ; the action \(x \mapsto \gamma_x\) gives the additive group \(A^+\) of \(A\) the structure of a group with operators with \(A\) as set of operators. The left ideals of \(A\) are just the subgroups of \(A^+\) which are stable under this action.

The left ideals in the ring A are just the right ideals of the opposite ring \(A^{0}\) . In a commutative ring the three species of ideals are the same; they are simply called ideals.

Examples. (1) Let A be a ring. The set A is a two-sided ideal of A; so is the set consisting of 0, which is called the zero ideal and sometimes denoted by 0 or (0) instead of \{0\}.

\begin{enumerate}
\def\labelenumi{(\arabic{enumi})}
\setcounter{enumi}{1}
\item
  For every element \(a\) of \(\mathbf{A}\) , the set \(\mathbf{A}.a\) of left multiples of \(a\) is a left ideal; similarly the set \(a.\mathbf{A}\) is a right ideal. When \(a\) is in the centre of \(\mathbf{A}\) , \(\mathbf{A}.a = a.\mathbf{A}\) ; this ideal is called the principal ideal generated by \(a\) and is denoted by (a). (a) = A if and only if \(a\) is invertible.
\item
  Let M be a subset of A. The set of elements \(x \in A\) such that xy = 0 for all \(y \in M\) is a left ideal of A called the left annihilator of M. The right annihilator of M is defined similarly.
\item
  Every intersection of left (resp. right, two-sided) of A is a left (resp. right, two-sided) ideal. Given a subset X of A, there thus exists a smallest left (resp. right, two-sided) ideal containing X; it is called the left (resp. right, two-sided) ideal generated by X.
\end{enumerate}

Let \(\mathfrak{a}\) be a left ideal of \(\mathbf{A}\) . The conditions \(1 \notin \mathfrak{a}, \mathfrak{a} \neq \mathbf{A}\) are obviously equivalent.

DEFINITION 5. Let A be a ring. By an abuse of language, a left ideal a is said to be maximal if it is a maximal element of the set of left ideals distinct from A.

In other words, a is maximal if \(a \neq A\) and the only left ideals of A containing a are a and A.

THEOREM 1 (Krull). Let A be a ring and a a left ideal of A distinct from A. There exists a maximal ideal m of A containing a.

Consider A as operating on the additive group \(A^{+}\) of A by left multiplication. Then the left ideals of A are the stable subgroups of \(A^{+}\) . The theorem thus follows from §4, no. 3, Proposition 3 applied to the subset \(P = \{1\}\) of \(A^{+}\) .

PROPOSITION 3. Let A be a ring, \((x_{\lambda})_{\lambda \in L}\) a family of elements of A and a (resp. b) the set of sums \(\sum_{\lambda \in L} a_{\lambda} x_{\lambda}\) where \((a_{\lambda})_{\lambda \in L}\) is a family with finite support of elements of A (resp. \(\sum_{\lambda \in L} a_{\lambda} x_{\lambda} b_{\lambda}\) where \((a_{\lambda})_{\lambda \in L}, (b_{\lambda})_{\lambda \in L}\) are families with finite support of elements of A). Then a (resp. b) is the left (resp. two-sided) ideal of A generated by the elements \(x_{\lambda}\) .

The formulae

\[
0 = \sum_ {\lambda \in L} 0. x _ {\lambda}\tag{18}
\]

\[
\sum_ {\lambda \in L} a _ {\lambda} x _ {\lambda} + \sum_ {\lambda \in L} a _ {\lambda} ^ {\prime} x _ {\lambda} = \sum_ {\lambda \in L} (a _ {\lambda} + a _ {\lambda} ^ {\prime}) x _ {\lambda}\tag{19}
\]

\[
a. \sum_ {\lambda \in L} a _ {\lambda} x _ {\lambda} = \sum_ {\lambda \in L} (a a _ {\lambda}) x _ {\lambda}\tag{20}
\]

prove that \(\mathfrak{a}\) is a left ideal. Let \(\mathfrak{a}'\) be a left ideal such that \(x_{\lambda} \in \mathfrak{a}'\) for all \(\lambda \in \mathbf{L}\) and let \((a_{\lambda})_{\lambda \in L}\) be a family with finite support in A. Then \(a_{\lambda}x_{\lambda}\in \mathfrak{a}'\) for all \(\lambda \in L\) , whence \(\sum_{\lambda \in L}a_{\lambda}x_{\lambda}\in \mathfrak{a}'\) ; hence \(\mathfrak{a}\subset \mathfrak{a}'\) . Hence \(\mathfrak{a}\) is the left ideal of A generated by the \(x_{\lambda}\) . The argument for \(\mathfrak{b}\) is analogous.

PROPOSITION 4. Let A be a ring and \((\mathfrak{a}_{\lambda})_{\lambda \in L}\) a family of left ideals of A. The left ideal generated by \(\bigcup_{\lambda \in L} \mathfrak{a}_{\lambda}\) consists of the sums \(\sum_{\lambda \in L} y_{\lambda}\) where \((y_{\lambda})_{\lambda \in L}\) is a family with finite support such that \(y_{\lambda} \in \mathfrak{a}_{\lambda}\) for all \(\lambda \in L\) .

Let \(\mathfrak{a}\) be the set of sums \(\sum_{\lambda \in \mathbf{L}} y_{\lambda}\) with \(y_{\lambda} \in \mathfrak{a}_{\lambda}\) for all \(\lambda \in \mathbf{L}\) . The formulae \(\sum_{\lambda \in \mathbf{L}} x_{\lambda} + \sum_{\lambda \in \mathbf{L}} y_{\lambda} = \sum_{\lambda \in \mathbf{L}} (x_{\lambda} + y_{\lambda})\) and \(a\) . \(\sum_{\lambda \in \mathbf{L}} x_{\lambda} = \sum_{\lambda \in \mathbf{L}} ax_{\lambda}\) shows that \(\mathfrak{a}\) is a left ideal of \(\mathbf{A}\) . Let \(\lambda \in \mathbf{L}\) and \(x \in \mathfrak{a}_{\lambda}\) ; write \(y_{\lambda} = x\) and \(y_{\mu} = 0\) for \(\mu \neq \lambda\) ; then \(x = \sum_{\lambda \in \mathbf{L}} y_{\lambda}\) , whence \(x \in \mathfrak{a}\) and finally \(\mathfrak{a}_{\lambda} \subset \mathfrak{a}\) . If a left ideal \(\mathfrak{a}'\) contains \(\mathfrak{a}_{\lambda}\) for all \(\lambda \in \mathbf{L}\) , it obviously contains \(\mathfrak{a}\) and hence \(\mathfrak{a}\) is generated by \(\bigcup_{\lambda \in \mathbf{L}} \mathfrak{a}_{\lambda}\) .

The ideal \(\mathfrak{a}\) generated by \(\bigcup_{\lambda \in L} \mathfrak{a}_{\lambda}\) is called the sum of the left ideals \(\mathfrak{a}_{\lambda}\) and is denoted by \(\sum_{\lambda \in L} \mathfrak{a}_{\lambda}\) (cf.~II, § 1, no. 7). In particular, the sum \(\mathfrak{a}_{1} + \mathfrak{a}_{2}\) of the two left ideals consists of the sums \(a_{1} + a_{2}\) where \(a_{1} \in \mathfrak{a}_{1}\) and \(a_{2} \in \mathfrak{a}_{2}\) .

\subsection{QUOTIENT RINGS}

Let A be a ring. If a is a two-sided ideal of A, two elements x and y of A are said to be congruent modulo a, written \(x \equiv y \pmod{a}\) or \(x \equiv y(a)\) , if \(x - y \in a\) . This is an equivalence relation on A. The relations \(x \equiv y(a)\) and \(x' \equiv y'(a)\) imply \(x + x' \equiv y + y'(a)\) , \(xx' \equiv xy'(a)\) for a is a left ideal and \(xy' \equiv yy'(a)\) for a is a right ideal, whence \(xx' \equiv yy'(a)\) . Conversely, if R is an equivalence relation on A compatible with addition and multiplication, the set a of \(x \equiv 0 \mod R\) is a two-sided ideal and \(x \equiv y \mod R\) is equivalent to \(x \equiv y \mod a\) .

Let A be a ring and a two-sided ideal of A. A/a denotes the quotient set of A by the equivalence relation \(x \equiv y(a)\) , with addition and multiplication the quotients of those on A (§ 1, no. 6, Definition 11). We show that A/a is a ring:

\begin{enumerate}
\def\labelenumi{(\alph{enumi})}
\item
  Under addition, A/a is the quotient commutative group of the additive group of A by the subgroup a.
\item
  Under multiplication, A/a is a monoid (§ 2, no. 1).
\item
  Let \(\xi, \eta, \zeta\) be in A/a and let \(\pi: \mathbf{A} \to \mathbf{A}/\mathfrak{a}\) be the canonical mapping; we choose elements \(x, y, z\) in A such that \(\pi(x) = \xi, \pi(y) = \eta\) and \(\pi(z) = \zeta\) . Then
\end{enumerate}

\[
\begin{array}{r l} \xi (\eta + \zeta) & = \pi (x) \pi (y + z) = \pi (x (y + z)) = \pi (x y + x z) \\ & = \pi (x) \pi (y) + \pi (x) \pi (z) = \xi \eta + \xi \zeta \end{array}
\]

and the relation \((\xi + \eta)\zeta = \xi \zeta + \eta \zeta\) is established similarly.

DEFINITION 6. Let A be a ring and a a two-sided ideal of A. The quotient ring of A by a, denoted by A/a, is the quotient set of A by the equivalence relation \(x \equiv y(a)\) , with addition and multiplication the quotients of those on A.

The ring A/\{0\} is isomorphic to A and A/A is a zero ring.

THEOREM 2. Let A be a ring and a two-sided ideal of A.

\begin{enumerate}
\def\labelenumi{(\alph{enumi})}
\item
  The canonical mapping \(\pi\) of A onto A/a is a ring homomorphism.
\item
  Let \(\mathbf{B}\) be a ring and \(f\) a homomorphism of \(\mathbf{A}\) into \(\mathbf{B}\) . If \(f(\mathfrak{a}) = \{0\}\) , there exists one and only one homomorphism \(\bar{f}\) of \(\mathbf{A} / \mathfrak{a}\) into \(\mathbf{B}\) such that \(f = \bar{f} \circ \pi\) .
\end{enumerate}

By construction, \(\pi(x + y) = \pi(x) + \pi(y)\) and \(\pi(xy) = \pi(x)\pi(y)\) for \(x, y\) in A; also \(\pi(1)\) is the unit \(\varepsilon\) of A/a, whence (a).

Let \(\mathbf{A}^{+}\) be the additive group of \(\mathbf{A}\) and \(\mathbf{B}^{+}\) that of \(\mathbf{B}\) ; as \(f\) is a homomorphism of \(\mathbf{A}^{+}\) into \(\mathbf{B}^{+}\) , zero on the subgroup \(a\) of \(\mathbf{A}^{+}\) , there exists (§ 4, no. 4, Proposition 5) one and only one homomorphism \(\bar{f}\) of \(\mathbf{A}^{+}/a\) into \(\mathbf{B}^{+}\) such that \(f = \bar{f} \circ \pi\) . Let \(\xi, \eta\) be in \(\mathbf{A}/a\) ; choose \(x, y\) in \(\mathbf{A}\) with \(\pi(x) = \xi\) and \(\pi(y) = \eta\) ; then \(\xi \eta = \pi(xy)\) , whence

\[
\bar {f} (\xi \eta) = \bar {f} (\pi (x y)) = f (x y) = f (x). f (y) = \bar {f} (\xi). \bar {f} (\eta)
\]

and \(\bar{f}(\varepsilon) = f(\pi(1)) = f(1)\) , hence \(\bar{f}\) is a ring homomorphism.

THEOREM 3. Let A and B be rings and f a homomorphism of A into B.

\begin{enumerate}
\def\labelenumi{(\alph{enumi})}
\item
  The kernel \(\mathfrak{a}\) off is a two-sided ideal of \(\mathbf{B}\) .
\item
  The image \(\mathbf{B}' = f(\mathbf{B})\) off is a subring of \(\mathbf{B}\) .
\item
  Let \(\pi: \mathbf{A} \to \mathbf{A} / \mathfrak{a}\) and \(i: \mathbf{B}' \to \mathbf{B}\) be the canonical morphisms. There exists one and only one morphism \(\bar{f}\) of \(\mathbf{A} / \mathfrak{a}\) into \(\mathbf{B}'\) such that \(f = i_{\circ} \bar{f} \circ \pi\) and \(\bar{f}\) is an isomorphism.
\end{enumerate}

As \(f\) is a morphism of the additive group of \(\mathbf{A}\) into that of \(\mathbf{B}\) , \(\mathfrak{a}\) is a subgroup of \(\mathbf{A}\) . If \(x \in \mathfrak{a}\) and \(a \in \mathbf{A}\) , then \(f(ax) = f(a)f(x) = 0\) , hence \(ax \in \mathfrak{a}\) and similarly \(xa \in \mathfrak{a}\) ; hence \(\mathfrak{a}\) is a two-sided ideal of \(\mathbf{A}\) . Assertion (b) is obvious. As \(f\) is zero on \(\mathfrak{a}\) , there exists a morphism \(\bar{f}\) of \(\mathbf{A} / \mathfrak{a}\) into \(\mathbf{B}'\) such that \(f = i \circ \bar{f} \circ \pi\) (Theorem 2). The uniqueness of \(\bar{f}\) and the fact that \(\bar{f}\) is an isomorphism follow from Set Theory, II, § 6, no. 4.

\subsection{SUBRINGS AND IDEALS IN A QUOTIENT RING}

PROPOSITION 5. Let A and A' be two rings, f a homomorphism of A into A' and a the kernel of f.

\begin{enumerate}
\def\labelenumi{(\alph{enumi})}
\item
  Let \(\mathbf{B}'\) be a subring of \(\mathbf{A}'\) . Then \(\mathbf{B} = f(\mathbf{B}')\) is a subring of \(\mathbf{A}\) containing \(\mathfrak{a}\) . If \(f\) is surjective, then \(f(\mathbf{B}) = \mathbf{B}'\) and \(f \mid \mathbf{B}\) defines when passing to the quotient an isomorphism of \(\mathbf{B}/\mathfrak{a}\) onto \(\mathbf{B}'\) .
\item
  Let \(\mathfrak{b}'\) be a left (resp. right, two-sided) ideal of \(\mathbf{A}'\) . Then \(\mathfrak{b} = f(\mathfrak{b}')\) is a left (resp. right, two-sided) ideal of \(\mathbf{A}\) containing \(\mathfrak{a}\) .
\item
  If \(\mathfrak{b}'\) is a two-sided ideal of \(\mathbf{A}'\) , the composite mapping of the canonical morphism \(\mathbf{A}' \to \mathbf{A}' / \mathfrak{b}'\) and \(f: \mathbf{A} \to \mathbf{A}'\) defines, when passing to the quotient, an injective morphism \(\bar{f}\) of \(\mathbf{A} / \mathfrak{b}\) into \(\mathbf{A}' / \mathfrak{b}'\) . If \(f\) is surjective, \(\bar{f}\) is an isomorphism of \(\mathbf{A} / \mathfrak{b}\) onto \(\mathbf{A}' / \mathfrak{b}'\) .
\item
  Suppose \(f\) is surjective. Let \(\Phi\) be the set of subrings (resp. left ideals, right ideals, two-sided ideals) of A containing a. Let \(\Phi'\) be the set of subrings (resp. left ideals, right ideals, two-sided ideals) of \(A'\) . The mappings \(B \to f(B)\) and \(B' \mapsto f'(B')\) are inverse bijections of \(\Phi\) onto \(\Phi'\) and \(\Phi'\) onto \(\Phi\) .
\end{enumerate}

(a) and (b) are obvious, except the last assertion of (a) which follows from no. 7, Theorem 3.

The composite morphism \(g: \mathbf{A} \to \mathbf{A}' \to \mathbf{A}' / \mathfrak{b}'\) considered in (c) has kernel \(\mathfrak{b}\) and hence \(\bar{f}\) is an injective morphism of \(\mathbf{A} / \mathfrak{b}\) into \(\mathbf{A}' / \mathfrak{b}'\) (§ 8, no. 7, Theorem 3). If \(f\) is surjective, \(g\) is surjective and hence \(\bar{f}\) is surjective.

Suppose \(f\) is surjective. By the above, the mapping \(\theta: \mathbf{B}' \mapsto f(\mathbf{B}')\) is a mapping of \(\Phi'\) into \(\Phi\) . Clearly the mapping \(\eta: \mathbf{B} \mapsto f(\mathbf{B})\) is a mapping of \(\Phi\) into \(\Phi'\) . Then \(\theta \circ \eta = \operatorname{Id}_{\Phi}, \eta \circ \theta = \operatorname{Id}_{\Phi'}\) , whence (d).

Remark. In the above notation, \(\theta\) and \(\eta\) are ordered set isomorphisms ( \(\Phi\) and \(\Phi'\) being ordered by inclusion).

COROLLARY. Let A be a ring and a a two-sided ideal of A.

\begin{enumerate}
\def\labelenumi{(\alph{enumi})}
\item
  Every left (resp. right, two-sided) ideal of A/a can be written uniquely in the form b/a, where b is a left (resp. right, two-sided) ideal of A containing a.
\item
  If \(\mathfrak{b}\) is two-sided, the composite homomorphism \(\mathbf{A} \to \mathbf{A} / \mathfrak{a} \to (\mathbf{A} / \mathfrak{a}) / (\mathfrak{b} / \mathfrak{a})\) defines when passing to the quotient an isomorphism of \(\mathbf{A} / \mathfrak{b}\) onto \((\mathbf{A} / \mathfrak{a}) / (\mathfrak{b} / \mathfrak{a})\) .
\end{enumerate}

It suffices to apply Proposition 5 to the canonical morphism of A onto A/a.

\subsection{MULTIPLICATION OF IDEALS}

Let A be a ring and a and b two-sided ideals of A. The set of elements of the form \(x_{1}y_{1} + \cdots + x_{n}y_{n}\) with \(n \geqslant 0\) , \(x_{i} \in a\) and \(y_{i} \in b\) for \(1 \leqslant i \leqslant n\) , is obviously a two-sided ideal of A, which is denoted by ab and called the product of the two-sided ideals a and b. Under this multiplication, the set of two-sided ideals of A is a monoid with unit element the two-sided ideal A. If a, b, c are two-sided ideals of A, then \(a(b + c) = ab + ac\) , \((b + c)a = ba + ca\) . If A is commutative, multiplication of ideals is commutative.

\(\mathsf{ab}\subset \mathsf{aA}\subset \mathsf{a}\) and \(\mathsf{ab}\subset \mathsf{Ab}\subset \mathsf{b}\) , hence

\[
a b \subset a \cap b.\tag{21}
\]

PROPOSITION 6. Let \(a, b_{1}, \ldots, b_{n}\) be two-sided ideals of A. If \(A = a + b_{i}\) for all \(i\) , then \(A = a + b_{1}b_{2} \ldots b_{n} = a + (b_{1} \cap b_{2} \cap \cdots \cap b_{n})\) .

By (21) it suffices to prove that \(\mathbf{A} = \mathfrak{a} + \mathfrak{b}_1\mathfrak{b}_2\ldots \mathfrak{b}_n\) . By induction, it suffices to consider the case where \(n = 2\) . By hypothesis, there exist \(a, a' \in \mathfrak{a}, b_1 \in \mathfrak{b}_1, b_2 \in \mathfrak{b}_2\) such that \(1 = a + b_1 = a' + b_2\) . Then

\[
1 = a ^ {\prime} + (a + b _ {1}) b _ {2} = \left(a ^ {\prime} + a b _ {2}\right) + b _ {1} b _ {2} \in \mathfrak {a} + \mathfrak {b} _ {1} \mathfrak {b} _ {2},
\]

whence \(\mathbf{A} = \mathfrak{a} + \mathfrak{b}_1\mathfrak{b}_2\)

PROPOSITION 7. Let \(\mathfrak{b}_1, \ldots, \mathfrak{b}_n\) be two-sided ideals of A such that \(\mathfrak{b}_i + \mathfrak{b}_j = \mathbf{A}\) for \(i \neq j\) . Then \(\mathfrak{b}_1 \cap \mathfrak{b}_2 \cap \dots \cap \mathfrak{b}_n = \sum_{\sigma \in \mathfrak{E}_n} \mathfrak{b}_{\sigma(1)} \mathfrak{b}_{\sigma(2)} \ldots \mathfrak{b}_{\sigma(n)}\) . In particular, if A is commutative, \(\mathfrak{b}_1 \cap \mathfrak{b}_2 \cap \dots \cap \mathfrak{b}_n = \mathfrak{b}_1 \mathfrak{b}_2 \ldots \mathfrak{b}_n\) (cf.~Exercise 2).

Suppose first \(n = 2\) . There exist \(a_1 \in \mathfrak{b}_1\) , \(a_2 \in \mathfrak{b}_2\) such that \(a_1 + a_2 = 1\) . If \(x \in \mathfrak{b}_1 \cap \mathfrak{b}_2\) , then \(x = x(a_1 + a_2) = xa_1 + xa_2 \in \mathfrak{b}_2\mathfrak{b}_1 + \mathfrak{b}_1\mathfrak{b}_2\) . Hence \(\mathfrak{b}_1 \cap \mathfrak{b}_2 = \mathfrak{b}_1\mathfrak{b}_2 + \mathfrak{b}_2\mathfrak{b}_1\) .

Suppose now the equation of the proposition is established for all integers \(<n\) . By Proposition 6, \(\mathfrak{b}_{n} + (\mathfrak{b}_{1}\mathfrak{b}_{2}\ldots \mathfrak{b}_{n - 1}) = A\) and hence

\[
\begin{array}{l} \mathfrak {b} _ {1} \cap \mathfrak {b} _ {2} \cap \dots \cap \mathfrak {b} _ {n} = (\mathfrak {b} _ {1} \cap \mathfrak {b} _ {2} \cap \dots \cap \mathfrak {b} _ {n - 1}) \mathfrak {b} _ {n} + \mathfrak {b} _ {n} (\mathfrak {b} _ {1} \cap \mathfrak {b} _ {2} \cap \dots \cap \mathfrak {b} _ {n - 1}) \\ = \Big (\sum_ {\sigma \in \mathfrak {E} _ {n - 1}} \mathfrak {b} _ {\sigma (1)} \mathfrak {b} _ {\sigma (2)} \dots \mathfrak {b} _ {\sigma (n - 1)} \Big) \mathfrak {b} _ {n} \\ \qquad + \mathfrak {b} _ {n} \Big (\sum_ {\sigma \in \mathfrak {E} _ {n - 1}} \mathfrak {b} _ {\sigma (1)} \mathfrak {b} _ {\sigma (2)} \dots \mathfrak {b} _ {\sigma (n - 1)} \Big) \\ \subset \sum_ {\tau \in \mathfrak {E} _ {n}} \mathfrak {b} _ {\tau (1)} \mathfrak {b} _ {\tau (2)} \dots \mathfrak {b} _ {\tau (n)} \subset \mathfrak {b} _ {1} \cap \mathfrak {b} _ {2} \cap \dots \cap \mathfrak {b} _ {n}. \end{array}
\]

\subsection{PRODUCT OF RINGS}

Let \((\mathbf{A}_i)_{i\in I}\) be a family of rings. Let A be the product set \(\prod_{i\in I} \mathbf{A}_i\) . On A addition and multiplication are defined by the formulae

\[
(x _ {i}) + (y _ {i}) = (x _ {i} + y _ {i}), \quad (x _ {i}) (y _ {i}) = (x _ {i} y _ {i}).\tag{22}
\]

It is immediately verified that A is a ring called the product of the rings \(A_{i}\) , with zero the element \(0 = (0_{i})_{i \in I}\) , where \(0_{i}\) is the zero of \(A_{i}\) , and unit \(1 = (1_{i})_{i \in I}\) , where \(1_{i}\) is the unit of \(A_{i}\) . If the \(A_{i}\) are commutative, so is A. If \(C_{i}\) is the centre of \(A_{i}\) , the centre of A is \(\prod_{i \in I} C_{i}\) .

For all \(i \in \mathbf{I}\) , the projection \(\mathbf{pr}_i\) of \(\mathbf{A}\) onto \(\mathbf{A}_i\) is a ring homomorphism. If \(\mathbf{B}\) is a ring and \(f_i \colon \mathbf{B} \to \mathbf{A}_i\) is a family of homomorphisms, there exists a unique homomorphism \(f \colon \mathbf{B} \to \mathbf{A}\) such that \(f_i = \mathbf{pr}_i \circ f\) for all \(i \in \mathbf{I}\) ; it is given by \(f(b) = (f_i(b))_{i \in \mathbf{I}}\) .

For all \(i \in I\) , let \(a_i\) be a left ideal of \(A_i\) . Then \(a = \prod_{i \in I} a_i\) is a left ideal of \(A\) . There is an analogous result for right ideals, two-sided ideals and subrings.

Suppose that \(a_{i}\) is a two-sided ideal for all \(i \in I\) and let \(f_{i}\) denote the canonical mapping of \(A_{i}\) onto \(A_{i}/a_{i}\) . Then the mapping \(f: (x_{i})_{i \in I} \mapsto (f_{i}(x_{i}))_{i \in I}\) of \(\prod_{i \in I} A_{i}\) onto \(\prod_{i \in I} (A_{i}/a_{i})\) is a ring homomorphism of kernel \(\prod_{i \in I} a_{i}\) and hence defines when passing to the quotient an isomorphism of \(\left( \prod_{i \in I} A_{i} \right) / \left( \prod_{i \in I} a_{i} \right)\) onto \(\prod_{i \in I} (A_{i}/a_{i})\) .

Let \((\mathbf{I}_{\lambda})_{\lambda \in \mathbb{L}}\) be a partition of I. The canonical bijection of \(\prod_{i\in I}\mathbf{A}_i\) onto \(\prod_{\lambda \in L}\left(\prod_{i\in \mathbf{I}_{\lambda}}\mathbf{A}_i\right)\) is a ring isomorphism, under which these two rings are identified. Let \(J\subset I\) . Let \(e_J\) denote the element \((x_i)_{i\in I}\) of A defined by \(x_{i} = 1_{i}\) for \(i\in J\) , \(x_{i} = 0_{i}\) for \(i\in I - J\) . Then \(e_J\) is a central idempotent (§ 1, no. 4) of A. The following formulae follow immediately:

\[
\begin{array}{r l} e _ {\mathrm{I}} = 1; \\ e _ {\varnothing} = 0; \\ e _ {\mathrm{J} \cap \mathrm{K}} = e _ {\mathrm{J}} e _ {\mathrm{K}} & \text { for   } \mathrm{J} \subset \mathrm{I}, \mathrm{K} \subset \mathrm{I}; \\ e _ {\mathrm{J} \cup \mathrm{K}} = e _ {\mathrm{J}} + e _ {\mathrm{K}} & \text { for   } \mathrm{J} \subset \mathrm{I}, \mathrm{K} \subset \mathrm{I}, \mathrm{J} \cap \mathrm{K} = \varnothing ; \\ \sum_ {\lambda} e _ {\mathrm{J} \lambda} = 1 & \text { if   } (\mathrm{J} _ {\lambda}) \text {   is   a   finite   partition   of   I }. \end{array}
\]

Let \(\mathbf{A}_{\mathbf{J}} = \prod_{i\in J}\mathbf{A}_i\) . Let \(\eta_{j}\) be the canonical projection of A onto \(\mathbf{A}_{\mathbf{J}}\) . For \(x = (x_{i})_{i\in J}\in \mathbf{A}_{\mathbf{J}}\) , let \(\varepsilon_{J}(x)\) be the element \((y_{i})_{i\in I}\) of A defined by \(y_{i} = x_{i}\) for \(i\in J, y_{i} = 0_{i}\) for \(i\in I - J\) . Then \(\eta_{j}\) is a ring homomorphism of A onto \(\mathbf{A}_{\mathbf{J}},\varepsilon_{j}\) is an injective homomorphism of the additive group of \(\mathbf{A}_{\mathbf{J}}\) onto A and in the diagram

\[
\mathrm{A} _ {\mathrm{J}} \xrightarrow {\varepsilon_ {\mathrm{J}}} \mathrm{A} \xrightarrow {\eta_ {\mathrm{I} - \mathrm{J}}} \mathrm{A} _ {\mathrm{I} - \mathrm{J}}
\]

the kernel \(a_{J}\) of \(\eta_{I-J}\) is equal to the image of \(\varepsilon_{J}\) . Then \(\varepsilon_{J}(xx') = \varepsilon_{J}(x)\varepsilon_{J}(x')\) for all \(x, x' \in A_{J}\) ; but \(\varepsilon_{J}\) is not in general a ring homomorphism for \(\varepsilon_{J}(1) = e_{J}\) . Clearly \(a_{J} = e_{J}A = Ae_{J}\) .

Let \(e_{\{i\}} = e_i\) and \(a_i = a_{\{i\}} = e_i A = Ae_i\) for all \(i \in I\) . Then \(e_i^2 = e_i\) ,

\[
e _ {i} e _ {j} = e _ {j} e _ {i} = 0
\]

for \(i \neq j\) . If I is finite, then \(\sum_{i \in I} e_i = 1\) , the additive group A is the direct sum of the two-sided ideals \(a_i\) and if \(x \in A\) its component in \(a_i\) is \(xe_i\) . The following proposition is immediately deduced.

PROPOSITION 8. Suppose I is finite. If b is a left or right ideal of A, b is the direct sum of the b ∩ a.

\subsection{DIRECT DECOMPOSITION OF A RING}

Let A be a ring and \((b_{i})_{i\in I}\) a family of two-sided ideals of A. We shall call the homomorphism

\[
x \mapsto (\phi_ {i} (x)) _ {i \in I},
\]

where \(\phi_{i}\) is the canonical homomorphism of A onto \(\mathbf{A} / \mathfrak{b}_i\) , the canonical homomorphism of A into \(\prod_{i\in I}(\mathbf{A} / \mathfrak{b}_i)\) .

PROPOSITION 9. Let A be a ring and \((\mathfrak{b}_1, \ldots, \mathfrak{b}_n)\) two-sided ideals of A such that \(\mathfrak{b}_i + \mathfrak{b}_j = \mathbf{A}\) for \(i \neq j\) . The canonical homomorphism of A into \(\prod_{i=1}^{n} (\mathbf{A} / \mathfrak{b}_i)\) is surjective of kernel \(\bigcap_{i=1}^{n} \mathfrak{b}_i = \sum_{\sigma \in \mathfrak{E}_n} \mathfrak{b}_{\sigma(1)} \mathfrak{b}_{\sigma(2)} \ldots \mathfrak{b}_{\sigma(n)}\) .

Clearly the kernel is \(\bigcap_{i=1}^{n} \mathfrak{b}_i\) . To prove the surjectivity, it is necessary to show that, for every family \((x_i)_{1 \leqslant i \leqslant n}\) of elements of A, there exists \(x \in \mathbf{A}\) such that \(x \equiv x_i (\mathfrak{b}_i)\) for all \(1 \leqslant i \leqslant n\) . We prove this assertion by induction on the cardinal \(n\) of I, the case \(n \leqslant 1\) being trivial. By the induction hypothesis, there exists \(y \in \mathbf{A}\) such that \(y \equiv x_i (\mathfrak{b}_i)\) for \(1 \leqslant i \leqslant n-1\) . We look for an \(x\) of the

form \(y + z\) with \(z \in \mathbf{A}\) . Of necessity \(z \equiv 0(\mathfrak{b}_i)\) for \(i < n\) , that is \(z \in \mathfrak{b} = \bigcap_{i=1}^{\infty} \mathfrak{b}_i\) , and on the other hand \(z \equiv x_n - y(\mathfrak{b}_n)\) . Now \(\mathfrak{b}_n + \mathfrak{b} = A\) by no. 9, Proposition 6, whence the existence of \(z\) . Finally, the second expression of the kernel follows from no. 9, Proposition 7.

DEFINITION 7. Let A be a ring. A finite family \((\mathfrak{b}_i)_{i\in I}\) of two-sided ideals of A such that the canonical homomorphism of A into \(\prod_{i\in I}(\mathbf{A} / \mathfrak{b}_i)\) is an isomorphism is called a direct decomposition of A.

PROPOSITION 10. Let A be a ring, A' its centre and \((\mathfrak{b}_{i})_{i\in I}\) a finite family of two-sided ideals of A. The following conditions are equivalent:

\begin{enumerate}
\def\labelenumi{(\alph{enumi})}
\item
  the family \((\mathfrak{b}_i)_{i\in I}\) is a direct decomposition of A;
\item
  there exists a family \((e_i)_{i\in I}\) of idempotents of \(\mathbf{A}'\) such that \(e_i e_j = 0\) for \(i\neq j\) , \(1 = \sum_{i\in I}e_i\) and \(\mathfrak{b}_i = \mathbf{A}(1 - e_i)\) for \(i\in I\) ;
\item
  \(\mathfrak{b}_i + \mathfrak{b}_j = \mathbf{A}\) for \(i \neq j\) and \(\bigcap_{i \in I} \mathfrak{b}_i = \{0\}\) ;
\item
  \(\mathfrak{b}_i + \mathfrak{b}_j = \mathbf{A}\) for \(i \neq j\) and \(\prod_{i \in I} \mathfrak{b}_i = \{0\}\) for every total order on \(\mathbf{I}\) ;
\item
  there exists a direct decomposition \((\mathfrak{b}_i')_{i\in \mathbf{I}}\) of \(\mathbf{A}'\) such that \(\mathfrak{b}_i = \mathbf{Ab}_i'\) for \(i\in \mathbf{I}\) .
\end{enumerate}

(a) \(\Rightarrow\) (b). If condition (a) holds, A may be identified with the ring \(\prod_{i\in I}(\mathbf{A} / b_i)\)

and \(b_{i}\) with the kernel of \(pr_{i}\) . The existence of the \(e_{i}\) with the properties in (b) then follows from no. 10.

\begin{enumerate}
\def\labelenumi{(\alph{enumi})}
\setcounter{enumi}{1}
\tightlist
\item
  \(\Rightarrow\) (d). Suppose that the \(e_i\) exist with the properties in (b). For \(i \neq j\) , \(1 - e_i \in \mathfrak{b}_i\) , \(e_i = e_i(1 - e_j) \in \mathfrak{b}_j\) , hence \(1 \in \mathfrak{b}_i + \mathfrak{b}_j\) and \(\mathbf{A} = \mathfrak{b}_i + \mathfrak{b}_j\) . On the other hand, if I is given a total ordering and \((x_i)_{i \in I}\) is a family of elements of A, then, since the \(e_i\) are central,
\end{enumerate}

\[
\prod_ {i \in \mathrm{I}} x _ {i} (1 - e _ {i}) = \Bigl (\prod_ {i \in \mathrm{I}} x _ {i} \Bigr) \Bigl (\prod_ {i \in \mathrm{I}} (1 - e _ {i}) \Bigr) = \Bigl (\prod_ {i \in \mathrm{I}} x _ {i} \Bigr) \Bigl (1 - \prod_ {i \in \mathrm{I}} e _ {i} \Bigr) = 0
\]

hence \(\prod_{i\in I}b_i = \{0\}\) .

\begin{enumerate}
\def\labelenumi{(\alph{enumi})}
\setcounter{enumi}{3}
\item
  \(\Rightarrow\) (c). This follows from no. 9, Proposition 7.
\end{enumerate}

(c) \(\Rightarrow\) (a). This follows from Proposition 9.

Thus conditions (a), (b), (c) and (d) are equivalent. Suppose they hold. By \((\mathbf{b}) \Rightarrow (\mathbf{a})\) , the family of \(\mathfrak{b}_i' = \mathbf{A}'(1 - e_i)\) is a direct decomposition of \(\mathbf{A}'\) . Then \(\mathfrak{b}_i = \mathbf{A}(1 - e_i) = \mathbf{Ab}_i'\) for all \(i \in \mathbf{I}\) . Hence condition (e) holds.

Finally, suppose condition (e) holds. By (a) \(\Rightarrow\) (b), there exists a family \((e_i)_{i\in I}\) of idempotents of \(A'\) such that \(e_i e_j = 0\) for \(i \neq j\) , \(1 = \sum_{i\in I} e_i\) and \(\mathfrak{b}_i' = A'(1 - e_i)\) for \(i \in I\) . Then \(\mathfrak{b}_i = Abi' = A(1 - e_i)\) for \(i \in I\) , hence condition (b) holds.

Remark. Let A be a ring. Let \((a_{i})_{i\in I}\) be a finite family of subgroups of the additive group \(A^{+}\) of A such that \(A^{+}\) is the direct sum of the \(a_{i}\) . Suppose \(a_{i}a_{i}\subset ai\) for \(i\in I\) and \(a_{i}a_{j} = \{0\}\) for \(i\neq j\) . Then \(a_{i}\) is for all \(i\in I\) a two-sided ideal of A. With addition and multiplication induced by those on A, \(a_{i}\) is a ring with unit element the component of \(1\in A\) in \(a_{i}\) . If \(b_{i} = \sum_{j\neq i}a_{j}\) , clearly the \(b_{i}\) satisfy condition (c) of Proposition 10 and hence \((b_{i})_{i\in I}\) is a direct decomposition of A, which is said to be defined by \((a_{i})_{i\in I}\) .

\BourbakiRunInHeading{Example: Ideals and quotient rings of Z}

An ideal of \(\mathbf{Z}\) is an additive subgroup of \(\mathbf{Z}\) and hence of the form \(n\) . \(\mathbf{Z}\) with \(n \geqslant 0\) ; conversely, for every integer \(n \geqslant 0\) , the set \(n\) . \(\mathbf{Z}\) is an ideal, the principal ideal \((n)\) . Thus every ideal of \(\mathbf{Z}\) is principal and is represented uniquely in the form \(n\mathbf{Z}\) with \(n \geqslant 0\) . The ideal (1) is equal to \(\mathbf{Z}\) , the ideal (0) consists of 0 and the ideals distinct from \(\mathbf{Z}\) and \(\{0\}\) are therefore of the form \(n\mathbf{Z}\) with \(n > 1\) . If \(m \geqslant 1\) and \(n \geqslant 1\) , \(m\mathbf{Z} \supset n\mathbf{Z}\) if and only if \(n \in m\) . \(\mathbf{Z}\) , that is \(m\) divides \(n\) . Therefore, for the ideal \(n\mathbf{Z}\) to be maximal, it is necessary and sufficient that there exist no integer \(m > 1\) distinct from \(n\) and dividing \(n\) ; in other words, the maximal ideals of \(\mathbf{Z}\) are the ideals of the form \(p\mathbf{Z}\) where \(p\) is a prime number (§ 4, no. 10, Definition 16).

Let \(m\) and \(n\) be two integers \(\geqslant 1\) . The ideal \(m\mathbf{Z} + n\mathbf{Z}\) is principal, whence there is an integer \(d \geqslant 1\) characterized by \(d\mathbf{Z} = m\mathbf{Z} + n\mathbf{Z}\) ; for every integer \(r \geqslant 1\) , the relation ``r divides d'' is equivalent to \(rZ \supset dZ\) and hence to `` \(rZ \supset mZ\) and \(rZ \supset nZ\) '', that is to ``r divides m and n''. It is thus seen that the common divisors of m and n are the divisors of d and that d is the greatest of the divisors \(\geqslant 1\) common to m and n; d is called the greatest common divisor (abbreviated to g.c.d.) of m and n.~As \(dZ = mZ + nZ\) , there exist two integers x and y such that \(d = mx + ny\) . m and n are said to be relatively prime if their g.c.d. is equal to 1. It amounts to the same to assume that there exist integers x and y with \(mx + ny = 1\) .

The intersection of the ideals mZ and nZ is non-zero for it contains mn and hence is of the form rZ with \(r \geqslant 1\) . Arguing as above, it is seen that the multiples of r are the common multiples of m and n and that r is the least of the integers \(\geqslant 1\) which are common multiples of m and n; it is called the least common multiple (l.c.m.) of m and n.

The product of the ideals \(m\mathbf{Z}\) and \(n\mathbf{Z}\) is the set of \(\sum_{i=1}^{r} mx_i ny_i = mn\left(\sum_{i=1}^{r} x_i y_i\right)\) for \(x_1, \ldots, y_r \in \mathbf{Z}\) and hence is equal to \(mn\mathbf{Z}\) .

For every integer \(n \geqslant 1\) , the quotient ring Z/nZ is called the ring of integers modulo n; it has n elements, which are the classes modulo n of the integers 0, 1, 2, \(\ldots\) , n - 1. For n = 1, we obtain the zero ring.

PROPOSITION 11. Let \(n_1, \ldots, n_d\) be integers \(\geqslant 1\) which are relatively prime in pairs and \(n = n_1 \ldots n_r\) . The canonical homomorphism of \(\mathbf{Z}\) into the product ring \(\prod_{i=1}^{r} \mathbf{Z} / n_i \mathbf{Z}\) is surjective of kernel \(n\mathbf{Z}\) and defines a ring isomorphism of \(\mathbf{Z} / n\mathbf{Z}\) onto \(\prod_{i=1}^{r} \mathbf{Z} / n_i \mathbf{Z}\) .

Let \(a_{i} = n_{i}\mathbf{Z}\) for \(i = 1, \ldots, r\) . By hypothesis, \(a_{i} + a_{j} = \mathbf{Z}\) for \(i \neq j\) . The proposition then follows from Proposition 9.

The above results, as also those concerning decomposition into prime factors, will be generalized in Chapter VII, § 1, which is devoted to the study of principal ideal domains and in Commutative Algebra, Chapter VII, § 3, which is devoted to the study of factorial domains.
